% !TEX program = xelatex
% !TEX encoding = UTF-8 Unicode
% -*- coding=utf-8 -*-
%% Compiler: Please use XeLaTeX to compile this document.
%% -----------------------------------------------------------------------------
%% BoyogalaElegantBook —— 中文使用手册
%%
%% 本文件既是使用说明，也是排版能力的演示文稿：手册中出现的每一种环境、
%% 每一个选项，都可以在正文中找到对应的例子。
%%
%% 编译方式（务必在项目根目录，也就是含有 fonts/ 目录的目录下执行）：
%%     xelatex boyogalaelegantbook-cn.tex
%%     latexmk -xelatex boyogalaelegantbook-cn.tex
%% -----------------------------------------------------------------------------

% 文档类选项：
%   lang=cn          简体中文（本章所用）
%   scheme=chinese   章号显示为“第一章”而非“第 1 章”
%   10pt             正文字号，其余字号随文档级别自动缩放
\documentclass[lang=cn,10pt,scheme=chinese]{boyogalaelegantbook}

% -----------------------------------------------------------------------------
% 封面信息：以下各项均为可选，未定义的项不会出现在封面上
% -----------------------------------------------------------------------------
\title{BoyogalaElegantBook：优美的 \LaTeX\ 书籍模板}   % 主标题
\subtitle{Elegant\LaTeX\ 经典之作}                    % 副标题
\author{Ethan Deng \& Liam Huang \& syvshc \& sikouhjw \& Osbert Wang}  % 作者
\institute{Elegant\LaTeX\ Program}                     % 机构
\date{2026/10/22}                                        % 日期
\version{1.0}                                          % 版本号
\bioinfo{自定义}{信息}                                  % 自定义信息条：{标题}{内容}

% 封面底部的提示文字
\extrainfo{\textcolor{red}{\bfseries 注意：BoyogalaElegantBook 仅支持 XeLaTeX 编译，且需要在包含 fonts/ 目录的项目根目录下编译。}}

% 目录显示到三级（章、节、小节）；类文件的默认值为 1（只显示章）
\setcounter{tocdepth}{3}

% 封面图与徽标：封面图按 1280x1024 裁剪，徽标为正方形
\logo{boyogalad.png}
\cover{cover.jpg}

% -----------------------------------------------------------------------------
% 本文档自行需要的宏包与命令
% 类文件只加载排版所必需的宏包，其余按需引入——这也是本项目的精简原则
% -----------------------------------------------------------------------------

% array 宏包提供表格列的扩展能力；\ccr 命令用来在表格中画一个色块
\usepackage{array}
\newcommand{\ccr}[1]{{\color{#1}\rule{1cm}{1cm}}}

% 画图章节需要 TikZ：类文件出于精简考虑不加载它，由文档按需引入
\usepackage{tikz}

% cprotect 允许在章节标题等位置使用 \lstinline 等脆弱命令
\usepackage{cprotect}

% 修改标题页中间色带的颜色（默认取配色主题的 second 色）
\definecolor{customcolor}{RGB}{32,178,170}
\colorlet{coverlinecolor}{customcolor}

% 参考文献数据库；打印文献的位置见正文末尾的 \printbibliography
\addbibresource[location=local]{reference.bib} % 参考文献，不要删除

\begin{document}

% 封面
\maketitle

% 前置部分：使用小写罗马数字编号，不计入章号
\frontmatter

% 目录
\tableofcontents

% 正文部分：重新使用阿拉伯数字编号
\mainmatter

%% 第一章：项目背景——来源、与上游的差异、获取与编译方式
\chapter{BoyogalaElegantBook 项目介绍}

\begin{quotation}
  \textbf{\textcolor{red}{BoyogalaElegantBook 是 ElegantBook 的精简与现代化分支：只保留一种字体配置，去除冗余选项与多余依赖，编译更快，跨平台排版结果一致。}}
\end{quotation}

BoyogalaElegantBook 源自 \href{https://github.com/ElegantLaTeX/ElegantBook}{ElegantBook}（由 \href{https://github.com/EthanDeng}{Ethan Deng} 与 \href{https://github.com/Liam0205}{Liam Huang} 创立，后由 \href{https://github.com/EthanDeng}{Ethan Deng}、\href{https://github.com/syvshc}{乙醇} 与 \href{https://github.com/sikouhjw}{死抠} 维护）。本项目在其基础上做了如下改动：

\begin{itemize}
  \item \textbf{唯一字体配置}：所有西文、中文与数学字体随项目提供，集中在 \lstinline{fonts/} 目录，不再有 \lstinline{math}、\lstinline{chinesefont} 等字体选项，也不再依赖系统字体；
  \item \textbf{仅支持 XeLaTeX}：使用 \lstinline{fontspec} 与 \lstinline{unicode-math} 统一加载字体；
  \item \textbf{精简依赖}：移除了 \lstinline{tikz}、\lstinline{mwe}、\lstinline{lipsum}、\lstinline{makecell}、\lstinline{hologo} 等非必要宏包，装饰符号改用数学字体中的字形；
  \item \textbf{修复缺陷}：修复了 \lstinline{color=none} 时的未定义颜色错误；语言模式精简为简体中文、繁体中文与英文三种，从源头消除了其它语言分支中的未定义命令问题；
  \item \textbf{全部字体均为 SIL OFL 1.1}：可自由使用、修改与再分发。
\end{itemize}

本文介绍本模板的设置内容与基本使用方法。如果有其他问题、建议或者意见，欢迎在 GitHub 上提交 \href{https://github.com/boyogala/BoyogalaElegantBook/issues}{issues}。

\begin{itemize}
  \item 项目地址：\href{https://github.com/boyogala/BoyogalaElegantBook}{https://github.com/boyogala/BoyogalaElegantBook}
  \item 上游项目：\href{https://github.com/ElegantLaTeX/ElegantBook}{ElegantBook}
  \item 联系邮箱：\mailto{boyogala@qq.com}
\end{itemize}

\begin{figure}[htbp]
  \centering
  \includegraphics[width=0.8\textwidth]{boyogala.png}
  \caption{微信搜一搜“博優旮旯”}
  \label{fig:project}
\end{figure}

\section{模板安装与更新}

本模板采用\textbf{免安装}方式使用：下载整个项目目录即可使用，无需安装到 \TeX\ 发行版中。

\subsection{在线使用模板}

把项目上传到 \href{https://www.overleaf.com/}{Overleaf} 即可在线使用。使用 Overleaf 的好处是无需安装 \TeX~Live，可以随时随地访问自己的文件。更多关于 Overleaf 的介绍和使用，请参考 Overleaf 的\href{https://www.overleaf.com/learn}{官方文档}。

\subsection{本地使用}

从 GitHub 下载项目压缩包并解压，得到如下目录结构：

\begin{lstlisting}
BoyogalaElegantBook/
├── boyogalaelegantbook.cls   % 类文件
├── boyogalaelegantbook-cn.tex % 中文手册
├── boyogalaelegantbook-en.tex % 英文手册
├── fonts/                     % 字体文件（必须保留）
└── figure/, image/            % 封面与插图
\end{lstlisting}

在项目根目录下用 XeLaTeX 编译即可：

\begin{lstlisting}
xelatex boyogalaelegantbook-cn.tex
% 或者
latexmk -xelatex boyogalaelegantbook-cn.tex
\end{lstlisting}

\textbf{请务必在项目根目录下编译}：字体路径是相对路径 \lstinline{./fonts/}，在其它目录中编译将找不到字体文件。

\subsection{环境要求}

本模板测试环境为
\begin{enumerate}
  \item Win11 + \TeX~Live 2023 及以上；
  \item Ubuntu 22.04/24.04 + \TeX~Live 2023 及以上；
  \item macOS + Mac\TeX\ 2023 及以上。
\end{enumerate}

模板需要 \lstinline{fontspec}、\lstinline{xeCJK}、\lstinline{ctex}、\lstinline{unicode-math}、\lstinline{tcolorbox}、\lstinline{titlesec}、\lstinline{fancyhdr}、\lstinline{tocloft}、\lstinline{biblatex} 等宏包，均包含在 \TeX~Live 的完整安装中。若使用精简安装，请补装上述宏包。

升级宏包的方法：使用 cmd 或 terminal 运行 \lstinline{tlmgr update --all}，如果 tlmgr 需要更新，请使用 cmd 运行 \lstinline{tlmgr update --self}，如果更新过程中出现了中断，请改用 \lstinline{tlmgr update --self --all --reinstall-forcibly-removed} 更新。

更多的内容请参考 \href{https://tex.stackexchange.com/questions/55437/how-do-i-update-my-tex-distribution}{How do I update my \TeX\ distribution?}

\subsection{关于 \CTeX\ 套装}

由于宏包版本问题，本模板不支持 \CTeX\ 套装，请使用 \TeX~Live/Mac\TeX\ 的完整安装。

%% 第二章：逐项说明模板的各类选项与用法
\chapter{BoyogalaElegantBook 设置说明}

本模板基于基础的 book 文类，所以 book 的选项对于本模板也是有效的（纸张无效，因为模板有设备选项）。默认编码为 UTF-8，推荐使用 \TeX~Live 编译。

%% 2.1 语言选项：简体中文、繁体中文与英文的启用方式及名称对照
\section{语言模式}
本模板只提供三种语言环境：\textbf{简体中文}（\lstinline{lang=cn}）、
\textbf{繁体中文}（\lstinline{lang=cht}）与\textbf{英文}（\lstinline{lang=en}，默认）。
改变语言环境会改变图表标题的引导词（图、表），文章结构词（目录、参考文献等），
以及定理环境中的引导词（定理、引理等）。

\begin{lstlisting}
\documentclass[cn]{boyogalaelegantbook}  % 简体中文
\documentclass[cht]{boyogalaelegantbook} % 繁体中文
\documentclass[lang=cn]{boyogalaelegantbook}
\end{lstlisting}

\begin{table}[htbp]
  \centering
  \caption{三种语言环境的结构名称对照}
  \label{tab:langs}
  \begin{tabular}{llll}
    \toprule
    结构 & 简体中文 & 繁体中文 & 英文 \\
    \midrule
    目录 & 目录 & 目錄 & Contents \\
    图 / 表 & 图 / 表 & 圖 / 表 & Figure / Table \\
    参考文献 & 参考文献 & 參考文獻 & Bibliography \\
    附录 & 附录 & 附錄 & Appendix \\
    定理 / 定义 / 引理 & 定理 / 定义 / 引理 & 定理 / 定義 / 引理 & Theorem / Definition / Lemma \\
    证明 & 证明 & 證明 & Proof \\
    注记 / 练习 / 习题 & 注记 / 练习 / 习题 & 註 / 練習 / 習題 & Remark / Exercise / Exercise \\
    \bottomrule
  \end{tabular}
\end{table}

\begin{remark}
简体中文与繁体中文使用\textbf{完全相同}的字体（见第 \ref{chap:font} 章），
差别只体现在上表的名称上；因此在繁体模式下书写的繁体字会与简体中文使用同一套字形，
若需要真正的台湾字形，请自行在导言区用 \lstinline{\\setCJKmainfont} 指定字体。
\end{remark}

%% 2.2 设备选项：A4 纸面与平板尺寸两种版面
\section{设备选项}
最早我们在 ElegantNote 模板中加入了设备选项（\lstinline{device}），后来，我们觉得这个设备选项的设置可以应用到 BoyogalaElegantBook 中\footnote{不过因为 BoyogalaElegantBook 模板封面图片的存在，在修改页面设计时，需要对图片进行裁剪。}，而且 Book 一般内容比较多，如果在 iPad 上看无需切边，放大，那用户的阅读体验将会得到巨大提升。你可以使用下面的选项将版面设置为 iPad 设备模式\footnote{默认为 normal 模式，也即 A4 纸张大小。}
\begin{lstlisting}
\documentclass[pad]{boyogalaelegantbook} %or
\documentclass[device=pad]{boyogalaelegantbook}
\end{lstlisting}

%% 2.3 配色主题：5 套预设配色与自定义配色
\section{颜色主题}

本模板内置 5 组颜色主题，分别为 \textcolor{structure1}{\lstinline{green}}\footnote{为原先默认主题。}、\textcolor{structure2}{\lstinline{cyan}}、\textcolor{structure3}{\lstinline{blue}}（默认）、\textcolor{structure4}{\lstinline{gray}}、\textcolor{structure5}{\lstinline{black}}。另外还有一个自定义的选项  \lstinline{nocolor}。调用颜色主题 \lstinline{green} 的方法为 
\begin{lstlisting}
\documentclass[green]{boyogalaelegantbook} %or
\documentclass[color=green]{boyogalaelegantbook}
\end{lstlisting}


\begin{table}[htbp]
  \caption{BoyogalaElegantBook 模板中的颜色主题\label{tab:color thm}}
  \centering
  \begin{tabular}{ccccccc}
  \toprule
    & \textcolor{structure1}{green} 
    & \textcolor{structure2}{cyan} 
    & \textcolor{structure3}{blue}
    & \textcolor{structure4}{gray} 
    & \textcolor{structure5}{black} 
    & 主要使用的环境\\
  \midrule
    structure & \ccr{structure1}
    & \ccr{structure2}
    & \ccr{structure3} 
    & \ccr{structure4} 
    & \ccr{structure5} 
    & chapter \ section \ subsection \\
    main      & \ccr{main1}
    & \ccr{main2}
    & \ccr{main3}
    & \ccr{main4}
    & \ccr{main5}
    & definition \ exercise \ problem \\
    second    & \ccr{second1}
    & \ccr{second2}
    & \ccr{second3}
    & \ccr{second4}
    & \ccr{second5}
    & theorem \ lemma \ corollary\\
    third     & \ccr{third1}
    & \ccr{third2}
    & \ccr{third3}
    & \ccr{third4}
    & \ccr{third5}
    & proposition\\
  \bottomrule
  \end{tabular}
\end{table}

如果需要自定义颜色的话请选择 \lstinline{nocolor} 选项或者使用 \lstinline{color=none}，然后在导言区定义 structurecolor、main、second、third 颜色，具体方法如下：
\begin{lstlisting}[tabsize=4]
\definecolor{structurecolor}{RGB}{0,0,0}
\definecolor{main}{RGB}{70,70,70}    
\definecolor{second}{RGB}{115,45,2}    
\definecolor{third}{RGB}{0,80,80}
\end{lstlisting}

%% 2.4 封面：可自由组合的封面元素、封面图、徽标与自定义封面
\section{封面}

\subsection{封面个性化}

从 3.10 版本开始，封面更加弹性化，用户可以自行选择输出的内容，包括 \lstinline{\title} 在内的所有封面元素都可为空。目前封面的元素有

\begin{table}[htbp]
  \centering
  \caption{封面元素信息}
  \begin{tabular}{p{0.07\textwidth}p{0.15\textwidth}|p{0.07\textwidth}p{0.15\textwidth}|p{0.07\textwidth}p{0.15\textwidth}}
    \hline
    信息 & 命令 & 信息 & 命令 & 信息 & 命令 \\
    \hline
    标题 & \lstinline|\title| & 副标题 & \lstinline|\subtitle| & 作者 & \lstinline|\author| \\
    机构 & \lstinline|\institute| & 日期 &  \lstinline|\date| & 版本 & \lstinline|\version| \\
    箴言 & \lstinline|\extrainfo| & 封面图 & \lstinline|\cover| & 徽标 & \lstinline|\logo| \\
    \hline
  \end{tabular}
\end{table}

另外，额外增加一个 \lstinline{\bioinfo} 命令，有两个选项，分别是信息标题以及信息内容。比如需要显示{\kaishu User Name：111520}，则可以使用 
\begin{lstlisting}
\bioinfo{User Name}{115520}
\end{lstlisting}

封面中间位置的色块的颜色可以使用下面命令进行修改：
\begin{lstlisting}
\definecolor{customcolor}{RGB}{32,178,170}
\colorlet{coverlinecolor}{customcolor}
\end{lstlisting}

\subsection{封面图}

本模板使用的封面图片来源于 \href{https://pixabay.com/en/tea-time-poetry-coffee-reading-3240766/}{pixabay.com}\footnote{感谢 China\TeX\ 提供免费图源网站，另外还推荐 \href{https://www.pexels.com/}{pexels.com}。}，图片完全免费，可用于任何场景。封面图片的尺寸为 $1280 \times 1024$, 更换图片的时候请\textbf{严格}按照封面图片尺寸进行裁剪。推荐一个免费的在线图片裁剪网站 \href{https://www.fotor.com/cn}{fotor.com}。用户 QQ 群内有一些合适尺寸的封面，欢迎取用。

\subsection{徽标}

本文用到的 Logo 比例为 1:1，也即正方形图片，在更换图片的时候请选择合适的图片进行替换。

\subsection{自定义封面}

另外，如果使用自定义的封面，比如 Adobe illustrator 或者其他软件制作的 A4 PDF 文档，请把 \lstinline{\maketitle} 注释掉，然后借助 \lstinline{pdfpages} 宏包将自制封面插入即可。如果使用 \lstinline{titlepage} 环境，也是类似。如果需要 2.x 版本的封面，请参考 \href{https://github.com/EthanDeng/etitlepage}{etitlepage}。

%% 2.5 章标题样式：hang 单行与 display 双行
\section{章标标题}

本模板内置 2 套\textit{章标题显示风格}，包含 \lstinline{hang}（默认）与 \lstinline{display} 两种风格，区别在于章标题单行显示（\lstinline{hang}）与双行显示（\lstinline{display}），本说明使用了 \lstinline{hang}。调用方式为
\begin{lstlisting}
\documentclass[hang]{boyogalaelegantbook} %or
\documentclass[titlestyle=hang]{boyogalaelegantbook}
\end{lstlisting}

在章标题内，章节编号默认是以数字显示，也即{\kaishu 第 1 章}，{\kaishu 第 2 章}等等，如果想要把数字改为中文，可以使用
\begin{lstlisting}
\documentclass[chinese]{boyogalaelegantbook} %or
\documentclass[scheme=chinese]{boyogalaelegantbook}
\end{lstlisting}

%% 2.6 定理环境：fancy 与 simple 两种模式、计数器与自定义定理
\section{数学环境简介}

在我们这个模板中，我们定义了两种不同的定理模式 \lstinline{mode}，包括简单模式（\lstinline{simple}）和炫彩模式（\lstinline{fancy}），默认为 \lstinline{fancy} 模式，不同模式的选择为
\begin{lstlisting}
\documentclass[simple]{boyogalaelegantbook} %or
\documentclass[mode=simple]{boyogalaelegantbook}
\end{lstlisting}

本模板定义了四大类环境

\begin{itemize}
  \item \textit{定理类环境}，包含标题和内容两部分，全部定理类环境的编号均以章节编号。根据格式的不同分为 3 种
    \begin{itemize}
      \item \textcolor{main}{\textbf{definition}} 环境，颜色为 \textcolor{main}{main}；
      \item \textcolor{second}{\textbf{theorem、lemma、corollary、axiom、postulate}} 环境，颜色为 \textcolor{second} {second}；
      \item \textcolor{third}{\textbf{proposition}} 环境，颜色为 \textcolor{third}{third}。
    \end{itemize}
  \item \textit{示例类环境}，有 \textbf{example、problem、exercise} 环境（对应于例、例题、练习），自动编号，编号以章节为单位，其中 \textbf{exercise} 有提示符。
  \item \textit{提示类环境}，有 \textbf{note} 环境，特点是：无编号，有引导符。
  \item \textit{结论类环境}，有 \textbf{conclusion、assumption、property、remark、solution} 环境\footnote{本模板还添加了一个 \lstinline|result| 选项，用于隐藏 \lstinline{solution} 和 \lstinline{proof} 环境，默认为显示（\lstinline{result=answer}），隐藏使用 \lstinline{result=noanswer}。}，三者均以粗体的引导词为开头，和普通段落格式一致。
\end{itemize}

其中，定理类环境均有带星号的版本：\textcolor{main}{\textbf{definition*}}、\textcolor{second}{\textbf{theorem*}}、\textcolor{second}{\textbf{lemma*}}、\textcolor{second}{\textbf{corollary*}}、\textcolor{second}{\textbf{axiom*}}、\textcolor{second}{\textbf{postulate*}}、\textcolor{third}{\textbf{proposition*}}，带星号的定理类环境不会编号。

\subsection{定理类环境的使用}

\subsubsection{\texttt{fancy} 模式}

在 \lstinline{fancy} 模式下使用了 \lstinline{tcolorbox} 宏包来定制定理类环境，所以和普通的定理环境的使用有些许区别，有编号定理的使用方法如下：

\begin{lstlisting}
% 有名字，有标签
\begin{theorem}{theorem name}{label}
  这是一个有名字和标签的定理。
  用 \ref{thm:label} 来引用这个定理。
\end{theorem}
% 无名字，有标签
\begin{theorem}{}{label no name}
  这是一个没有名字但是有标签的定理。
  用 \ref{thm:label no name} 来引用这个定理。
\end{theorem}
% 有名字，无标签
\begin{theorem}{theorem name}{}
  这是一个有名字但是没有标签的定理。
  这个定理不能被引用。
  最后的 {} 可以省略不写。
\end{theorem}
% 无名字，无标签
\begin{theorem}{}{}
  这是一个没有名字也没有标签的定理。
  这个定理不能被引用。
  两个 {} 均可以省略不写。
\end{theorem}
\end{lstlisting}

第一个选项 \lstinline{theorem name} 是定理的名字，如果定理没有名字请使用 \lstinline|{||\cprotect\footnote{除非这个定理也没有标签，否则不能省略。}。第二个选项 \lstinline{label} 是交叉引用时所用到的标签，如果定理没有标签，可以不写 \lstinline|{label|| 或使用 \lstinline|{||。交叉引用的方法为 \lstinline|\ref{thm:label||。请注意，交叉引用时必须加上前缀 \lstinline{thm:}。其他相同用法的定理类环境见表 \ref{tab:theorem-class}。

不编号的定理使用方法如下：

\begin{lstlisting}
% 有名字
\begin{theorem*}{theorem name}
  这是一个不编号的有名字的定理。
  这个定理不能被引用。
\end{theorem*}
% 无名字
\begin{theorem*}{}
  这是一个不编号且没有名字的定理。
  最后的 {} 可以省略不写。
  这个定理不能被引用。
\end{theorem*}
\end{lstlisting}
其中的选项 \lstinline{theorem name} 是定理的名字，如果没有名字可以不写 \lstinline|{theorem name|| 或使用 \lstinline|{||

\begin{table}[htbp]
  \centering
  \caption{定理类环境}
    \begin{tabular}{llll}
    \toprule
    环境名 & 标签名 & 前缀 & 交叉引用 \\
    \midrule
    definition & label & def   & \lstinline|\ref{def:label|| \\
    theorem & label & thm   & \lstinline|\ref{thm:label|| \\
    postulate & label & pos & \lstinline|\ref{pos:label|| \\
    axiom & label & axi & \lstinline|\ref{axi:label||\\
    lemma & label & lem   & \lstinline|\ref{lem:label|| \\
    corollary & label & cor   & \lstinline|\ref{cor:label|| \\
    proposition & label & pro   & \lstinline|\ref{pro:label|| \\
    \bottomrule
    \end{tabular}%
  \label{tab:theorem-class}%
\end{table}%

在用户多次反馈下，在 4.1 版本引入了原生定理的支持方式，使用方法如下：

\begin{lstlisting}
% 有名字，有标签
\begin{theorem}[theorem name]\label{thm:label}
  这是一个有名字和标签的定理。
  用 \ref{thm:label} 来引用这个定理。
\end{theorem}
% 无名字，有标签
\begin{theorem}\label{thm:label no name}
  这是一个没有名字但是有标签的定理。
  用 \ref{thm:label no name} 来引用这个定理。
\end{theorem}
% 有名字，无标签
\begin{theorem}[theorem name]
  这是一个有名字但是没有标签的定理。
  这个定理不能被引用。
\end{theorem}
% 无名字，无标签
\begin{theorem}
  这是一个没有名字也没有标签的定理。
  这个定理不能被引用。
\end{theorem}
\end{lstlisting}
这时引用不需要加前缀，\lstinline{\ref} 中的内容与 \lstinline{\label} 中的内容相同，换句话说，这时的使用方法与使用原生定理环境的时候几乎相同。

\subsubsection{\texttt{simple} 模式}

在 \lstinline{simple} 模式下使用了 \lstinline{amsthm} 宏包，定理类环境使用方法与原生一致，定理的使用方法如下：

\begin{lstlisting}
% 有名字，有标签
\begin{theorem}[theorem name]\label{thm:label}
  这是一个有名字和标签的定理。
  用 \ref{thm:label} 来引用这个定理。
\end{theorem}
% 无名字，有标签
\begin{theorem}\label{thm:label no name}
  这是一个没有名字但是有标签的定理。
  用 \ref{thm:label no name} 来引用这个定理。
\end{theorem}
% 有名字，无标签
\begin{theorem}[theorem name]
  这是一个有名字但是没有标签的定理。
  这个定理不能被引用。
\end{theorem}
% 无名字，无标签
\begin{theorem}
  这是一个没有名字也没有标签的定理。
  这个定理不能被引用。
\end{theorem}
\end{lstlisting}

% \subsection{算法环境}

% \begin{algorithm}\label{alg:test}
%   \Input{A bitmap $I$ of size $w \times l$}
%   \Output{A partition of the bitmap}
%   \BlankLine
%   \emph{special treatment of the first line}\;
%   \For{$i \leftarrow 2$ \KwTo $l$}{
%     \emph{special treatment of the first element of line $i$}\;
%     \For{$j \leftarrow 2$ \KwTo $w$}{\label{forins}
%       $\Left \leftarrow \FindCompress{$I[i,j-1]$}$\;
%       $\Up \leftarrow \FindCompress{$I[i-1,]$}$\;
%       $\This \leftarrow \FindCompress{$I[i,j]$}$\;
%       \If(\tcp*[h]{O(\Left,\This)==1}){\Left compatible with \This}{\label{lt}
%         \lIf{$\Left < \This$}{$\Union{\Left,\This}$}
%         \lElse{$\Union{\This,\Left}$}
%       }
%       \If(\tcp*[f]{O(\Up,\This)==1}){\Up compatible with \This}{\label{ut}
%         \lIf{$\Up < \This$}{$\Union{\Up,\This}$}
%         \tcp{\This is put under \Up to keep tree as flat as possible}\label{cmt}
%         \lElse{$\Union{\This,\Up}$}\tcp*[r]{\This{} linked to \Up}\label{lelse}
%       }
%     }
%     \lForEach{element $e$ of the line $i$}{\FindCompress{p}}
%   }
%   \caption{disjoint decomposition}\label{algo_disjdecomp}
% \end{algorithm}


\subsection{修改计数器}

当前定理等环境计数器按章计数，如果想修改定理类环境按节计数，可以修改计数器选项 \lstinline{thmcnt}，可用选项为 \lstinline{chapter} （默认）与 \lstinline{section}：

\begin{lstlisting}
\documentclass[section]{boyogalaelegantbook} %or
\documentclass[thmcnt=section]{boyogalaelegantbook}
\end{lstlisting}

如果希望全局的定理类环境使用同一个计数器，可以使用文档类选项 \lstinline{usesamecnt}：

\begin{lstlisting}
  \documentclass[usesamecnt]{boyogalaelegantbook}
\end{lstlisting}

\subsection{自定义定理类环境}

4.4 版本新增了一个自定义定理类环境的命令：\lstinline|\elegantnewtheorem|，它的参数含义如下：

\begin{lstlisting}
% fancy 模式（默认）
\elegantnewtheorem{env}{title}{style}{prefix}[numbered-like]
% simple 模式
\elegantnewtheorem{env}{title}{style}[numbered-like]
\end{lstlisting}
该命令可以同时定义编号环境 \lstinline|env| 和不编号环境 \lstinline|env*|。

其中 \lstinline|style| 支持的参数有：\lstinline|defstyle|，\lstinline|thmstyle|，\lstinline|prostyle|，分别对应“定义”，“定理”，“命题”三种样式。

如果添加了可选参数 \lstinline{numbered-like}，将会使该定理类环境与名为 \lstinline{numbered-like} 的定理类环境使用同一计数器。\textbf{注意}：该参数在使用 \lstinline{usesamecnt} 选项时不起作用，并且会在终端以及 \lstinline{.log} 文件中输出一个警告，来提示用户该选项不起作用：

\begin{lstlisting}
  [numbered-like] won't make sence with option
  `usesamecnt'.
\end{lstlisting}

\begin{itemize}
  \item 在炫彩模式（\lstinline{fancy}）下，需要 5 个参数来定义一个新的定理类环境，分别是：
  定理类环境名，定理类环境的标题，定理类环境的样式，该定理类环境的前缀，（可选）该定理类环境继承的定理类环境：

\begin{lstlisting}
% 导言区
\elegantnewtheorem{examplefancy}{自定义定理类环境}{thmstyle}{exfancy}
% 正文
\begin{examplefancy}{定理名}{label}
  这里是自定义定理类环境 \ref{exfancy:label}
\end{examplefancy}
\begin{examplefancy*}{定理名}
  这里是无编号自定义定理类环境
\end{examplefancy*}
\end{lstlisting}

  如果不给出第四个参数，或第四个参数置空的话，将会用定理类环境名来作为默认前缀，即
\begin{lstlisting}
% 导言区
\elegantnewtheorem{test}{TEST}{thmstyle}
% 或
\elegantnewtheorem{test}{TEST}{thmstyle}{}
% 正文
\begin{test}{name}{label}
  默认前缀为 test。
  使用 \ref{test:label} 来引用这个定理类环境。
\end{test}
\end{lstlisting}

这时会在终端以及 \verb|.log| 文件中输出一个警告信息来提示用户没有定义前缀：

\begin{lstlisting}[language=bash]
Class elegantbook Warning: Because you didn't provide a prefix. 
(elegantbook)              We use test as the default prefix. 
(elegantbook)              You have to use 
(elegantbook)              \ref {test:label} to refer a 
(elegantbook)              \begin {test}{name}{label} environment. 
(elegantbook)               on input line 3.
\end{lstlisting}

  \item 在简单模式（\lstinline{simple}）下，需要 4 个参数来定义一个新的定理类环境，分别是：
  定理类环境名，定理类环境的标题，定理类环境的样式，该定理类环境的前缀，（可选）该定理类环境继承的定理类环境：
\begin{lstlisting}
% 导言区
\elegantnewtheorem{examplesimple}{自定义定理类环境}{thmstyle}
% 正文
\begin{examplesimple}[定理名]\label{exsimple:label}
  这里是自定义定理类环境 \ref{exsimple:label}
\end{examplesimple}
\begin{examplesimple*}[定理名]
  这里是无编号自定义定理类环境
\end{examplesimple*}
\end{lstlisting}

  如果此时错误地给出了第四个参数，那么将会在终端以及 \verb|.log| 文件中输出一个错误信息：
\begin{lstlisting}
% elegantbook-cn.tex
\elegantnewtheorem{test}{TEST}{thmstyle}{}
% .log file
./elegantbook-cn.tex:3: Class elegantbook Error: You can't set a prefix in mode ``simple''.
(elegantbook)              Just use 
(elegantbook)              \elegantnewtheorem {test}{TEST}{thmstyle} .
\end{lstlisting}

\end{itemize}

\subsection{其他环境的使用}

其他三种环境没有选项，可以直接使用，比如 \lstinline{example} 环境的使用方法与效果：
\begin{lstlisting}
\begin{example}
   This is the content of example environment.
\end{example}
\end{lstlisting}

这几个都是同一类环境，区别在于

\begin{itemize}
  \item 示例环境（example）、练习（exercise）与例题（problem）章节自动编号；
  \item 注意（note），练习（exercise）环境有提醒引导符；
  \item 结论（conclusion）等环境都是普通段落环境，引导词加粗。
\end{itemize}

%% 2.7 列表环境：itemize 与 enumerate 的层次与符号
\section{列表环境}
本模板借助于 \lstinline{tikz} 定制了 \lstinline{itemize} 和 \lstinline{enumerate} 环境，其中 \lstinline{itemize} 环境修改了 3 层嵌套，而 \lstinline{enumerate} 环境修改了 4 层嵌套（仅改变颜色）。示例如下\\[2ex]
\begin{minipage}[b]{0.49\textwidth}
  \begin{itemize}
    \item first item of nesti;
    \item second item of nesti;
      \begin{itemize}
        \item first item of nestii;
        \item second item of nestii;
        \begin{itemize}
          \item first item of nestiii;
          \item second item of nestiii.
        \end{itemize}   
      \end{itemize}
  \end{itemize}
\end{minipage}
\begin{minipage}[b]{0.49\textwidth}
  \begin{enumerate}
    \item first item of nesti;
    \item second item of nesti;
      \begin{enumerate}
        \item first item of nestii;
        \item second item of nestii;
        \begin{enumerate}
          \item first item of nestiii;
          \item second item of nestiii.
        \end{enumerate}   
      \end{enumerate}
  \end{enumerate}
\end{minipage}

%% 2.8 参考文献：biblatex 的打印方式与样式选项
\section{参考文献}

\subsection{打印文献}

之前我们将文献调用的命令放在模板里面，然后用户反馈 \lstinline{\cite} 命令无法自动补全，因此我们新版本将其拿到外面来，新版本打印参考文献的命令的方法是，在导言区（也即 \lstinline|\begin{document|| 之前），加入：

\begin{lstlisting}
  \addbibresource[location=local]{reference.bib}
\end{lstlisting}

然后再需要打印文献的地方使用：
\begin{lstlisting}
  \printbibliography[heading=bibintoc, title=\ebibname]
\end{lstlisting}

其中 \lstinline{reference.bib} 为参考文献存放的文件，需要放在项目文件夹下。

\subsection{修改文献格式}

此外，本模板调用了 biblatex 宏包，并提供了 biber（默认） 和 bibtex 两个后端选项，可以使用 \lstinline{bibend} 进行修改：

\begin{lstlisting}
  \documentclass[bibtex]{boyogalaelegantbook}
  \documentclass[bibend=bibtex]{boyogalaelegantbook}
\end{lstlisting}

关于文献条目（bib item），你可以在谷歌学术，Mendeley，Endnote 中取，然后把它们添加到 \lstinline{reference.bib} 中。在文中引用的时候，引用它们的键值（bib key）即可。

为了方便文献样式修改，模板引入了 \lstinline{bibstyle} 和 \lstinline{citestyle} 选项，默认均为数字格式（numeric），参考文献示例：\cite{cn1,en2,en3} 使用了中国一个大型的 P2P 平台（人人贷）的数据来检验男性投资者和女性投资者在投资表现上是否有显著差异。

如果需要设置为国标 GB7714-2015，需要使用：
\begin{lstlisting}
  \documentclass[citestyle=gb7714-2015, bibstyle=gb7714-2015]{boyogalaelegantbook} 
\end{lstlisting}

在使用
\begin{lstlisting}
  \documentclass[citestyle=gb7714-2015, bibstyle=gb7714-2015]{boyogalaelegantbook} 
\end{lstlisting}
后，排序方式为按引用先后排序。如果不使用国标 GB7714-2015 并且需要添加排序方式，可以在导言区加入下面命令：
\begin{lstlisting}
  \ExecuteBibliographyOptions{sorting=<name>}
\end{lstlisting}
其中 \lstinline{<name>} 可以是 nty, nyt, nyvt, anyt, anyvt, ynt, ydnt, none, count, debug 的其中之一，更多排序相关请参考 biblatex 宏包手册的 3.1.2.1 节。


%% 2.9 序章：加入不改变章节编号的前言
\section{添加序章}

如果你想在第一章前面添序章，不改变原本章节序号，可以在第一章内容前面使用 
\begin{lstlisting}
\chapter*{Introduction}
\markboth{Introduction}{Introduction}
The content of introduction.
\end{lstlisting}

%% 2.10 目录：单栏与双栏，以及目录深度
\section{目录选项与深度}
本模板添加了一个目录选项 \lstinline{toc}，可以设置目录为单栏（\lstinline{onecol}）和双栏（\lstinline{twocol}）显示，比如双栏显示可以使用
\begin{lstlisting}
\documentclass[twocol]{boyogalaelegantbook}
\documentclass[toc=twocol]{boyogalaelegantbook}
\end{lstlisting}

默认本模板目录深度为 1，你可以在导言区使用
\begin{lstlisting}
\setcounter{tocdepth}{2}
\end{lstlisting}
将其修改为 2 级目录（章与节）显示。


%% 2.11 章节摘要：introduction 环境
\section{章节摘要}
模板新增了一个章节摘要环境（introduction），使用示例
\begin{lstlisting}
\begin{introduction}
  \item Definition of Theorem
  \item Ask for help
  \item Optimization Problem
  \item Property of Cauchy Series
  \item Angle of Corner
\end{introduction}
\end{lstlisting}
效果如下：
\begin{introduction}
  \item Definition of Theorem
  \item Ask for help
  \item Optimization Problem
  \item Property of Cauchy Series
  \item Angle of Corner
\end{introduction}

环境的标题文字可以通过这个环境的可选参数进行修改，修改方法为：
\begin{lstlisting}
\begin{introduction}[Brief Introduction]
...
\end{introduction}
\end{lstlisting}

%% 2.12 章后习题：problemset 环境
\section{章后习题}
前面我们介绍了例题和练习两个环境，这里我们再加一个，章后习题（\lstinline{problemset}）环境，用于在每一章结尾，显示本章的练习。使用方法如下

\begin{lstlisting}
\begin{problemset}
  \item exercise 1
  \item exercise 2
  \item exercise 3
\end{problemset}
\end{lstlisting}


效果如下：
\begin{problemset}[我的题目]
  \item exercise 1
  \item exercise 2
  \item exercise 3
  \item 测试数学公式
  \begin{equation}
    a^2+b^2=c_{2_{i}} (1,2) [1,23]
  \end{equation}
\end{problemset}

\begin{remark}
如果你想把 \lstinline{problemset} 环境的标题改为其他文字，你可以类似于 introduction 环境修改 problemset 的可选参数。另外，目前这个环境会自动出现在目录中，但是不会出现在页眉页脚信息中（待解决）。
\end{remark}

\begin{solution}
如果你想把 \lstinline{problemset} 环境的标题改为其他文字，你可以类似于 introduction 环境修改 problemset 的可选参数。另外，目前这个环境会自动出现在目录中，但是不会出现在页眉页脚信息中（待解决）。
\end{solution}

%% 2.13 旁注：右侧边注区的开启与使用
\section{旁注}

在 3.08 版本中，我们引入了 旁注设置选项 \lstinline{marginpar=margintrue} 以及测试命令 \lstinline{\elegantpar} ，但是由此带来一堆问题。我们决定在 3.09 版本中将其删除，并且，在旁注命令得到大幅度优化之前，不会将此命令再次引入书籍模板中。对此造成各位用户的不方便，非常抱歉！不过我们保留了 \lstinline{marginpar} 这个选项，你可以使用 \lstinline{marginpar=margintrue} 获得保留右侧旁注的版面设计。然后使用系统自带的 \lstinline{\marginpar} 或者 \lstinline{marginnote} 宏包的 \lstinline{\marginnote} 命令。

\begin{remark}
在使用旁注的时候，需要注意的是，文本和公式可以直接在旁注中使用。

\begin{lstlisting}
% text
\marginpar{margin paragraph text}

% equation
\marginpar{
  \begin{equation}
    a^2 + b^2 = c^2
  \end{equation}
}
\end{lstlisting}

但是浮动体（表格、图片）需要注意，不能用浮动体环境，需要使用直接插图命令或者表格命令环境。然后使用 \lstinline{\captionof} 为其设置标题。为了得到居中的图表，可以使用 \lstinline{\centerline} 命令或者 \lstinline{center} 环境。更多详情请参考：\href{https://tex.stackexchange.com/questions/5583/caption-of-figure-in-marginpar-and-caption-of-wrapfigure-in-margin}{Caption of Figure in Marginpar}。

\begin{lstlisting}
% graph with centerline command
\marginpar{
  \centerline{
    \includegraphics[width=0.2\textwidth]{logo.png}
  }
  \captionof{figure}{your figure caption}
}

% graph with center environment
\marginpar{
  \begin{center}
    \includegraphics[width=0.2\textwidth]{logo.png}
    \captionof{figure}{your figure caption}
  \end{center}
}
\end{lstlisting}

\end{remark}

%% 第三章：唯一字体配置——字体清单、编译引擎、中文四体与数学字体
\chapter{字体方案}
\label{chap:font}
BoyogalaElegantBook 只保留\textbf{一种}字体配置：所有西文、中文与数学字体都随模板一起提供，放在项目根目录下的 \lstinline{fonts/} 目录中。编译时模板直接从该目录读取字体文件，因此不再依赖操作系统里安装的字体，在任何平台上的排版结果都完全一致。

\textcolor{red}{\bfseries 重要提示}：由于字体全部来自 \lstinline{fonts/} 目录，\textbf{请务必在项目根目录（也就是 \lstinline{fonts/} 所在目录）下编译文档}，否则模板将找不到字体文件。

\section{字体清单}

所有字体均以 SIL Open Font License 1.1 发布，可以自由使用、修改与再分发；版权声明与许可证全文见 \lstinline{fonts/OFL.txt}。

\begin{table}[htbp]
  \centering
  \caption{BoyogalaElegantBook 的字体清单}
  \begin{tabular}{llll}
  \toprule
  用途 & 字体 & 文件名 & 授权 \\
  \midrule
  西文正体 & STIX Two Text & \lstinline|STIXTwoText-*.otf| & OFL 1.1 \\
  西文无衬线 & Source Sans 3 & \lstinline|SourceSans3-*.otf| & OFL 1.1 \\
  等宽字体 & JetBrains Mono & \lstinline|JetBrainsMono-*.ttf| & OFL 1.1 \\
  数学字体 & STIX Two Math & \lstinline|STIXTwoMath-Regular.otf| & OFL 1.1 \\
  中文宋体 & 思源宋体 & \lstinline|SourceHanSerifSC-Regular.otf| & OFL 1.1 \\
  中文黑体 & 思源黑体 & \lstinline|SourceHanSansSC-*.otf| & OFL 1.1 \\
  中文楷体 & 霞鹜文楷 & \lstinline|LXGWWenKai-Regular.ttf| & OFL 1.1 \\
  中文仿宋 & 朱雀仿宋 & \lstinline|ZhuqueFangsong-Regular.ttf| & OFL 1.1 \\
  \bottomrule
  \end{tabular}
\end{table}

\section{编译引擎}

本模板\textbf{只支持 XeLaTeX}。所有字体都通过 \lstinline{fontspec} 与 \lstinline{unicode-math} 加载，使用其它引擎编译时模板会直接报错并提示改用 XeLaTeX。

\begin{lstlisting}
xelatex boyogalaelegantbook-cn.tex
% 或者
latexmk -xelatex boyogalaelegantbook-cn.tex
\end{lstlisting}

\section{中文字体的四种角色}

中文仍然保留宋、黑、楷、仿宋四种角色，与传统书籍的排版习惯一致：

\begin{itemize}
  \item \textbf{宋体}：正文；
  \item \textbf{黑体}：粗体与标题；
  \item \textbf{楷体}：定理类环境正文，以及斜体；
  \item \textbf{仿宋}：证明环境正文。
\end{itemize}

对应的命令为 \lstinline|\songti|、\lstinline|\heiti|、\lstinline|\kaishu|、\lstinline|\fangsong|，它们在 \lstinline|lang=cn| 与 \lstinline|lang=cht| 模式下可用：

\begin{lstlisting}
{\songti 这是宋体} \quad {\heiti 这是黑体} \quad
{\kaishu 这是楷体} \quad {\fangsong 这是仿宋}
\end{lstlisting}

\section{数学字体}

数学字体使用 STIX Two Math，由 \lstinline{unicode-math} 加载，提供了完整的 AMS 符号集以及黑体、哥特体、花体、脚本体等数学字母：

\begin{equation}
  \mathbb{R} \subset \mathcal{C}, \qquad
  \mathfrak{g} \in \mathscr{L}, \qquad
  \bm{x} \in \mathbb{R}^n.
\end{equation}

由于 \lstinline{bm} 宏包与 \lstinline{unicode-math} 不兼容，模板不再加载 \lstinline{bm| 宏包，而是把 \lstinline|\bm| 定义为 \lstinline{unicode-math} 提供的 \lstinline|\symbf|，用法与原来完全一致。

\section{自定义字体}

如果需要替换字体，可以在导言区重新设置。例如改用系统安装的字体：

\begin{lstlisting}
\setmainfont{TeX Gyre Termes}[
  Extension = .otf,
  UprightFont = *-regular,
  BoldFont = *-bold,
  ItalicFont = *-italic,
  BoldItalicFont = *-bolditalic]
\setCJKmainfont[
  Path = ./fonts/,
  BoldFont = SourceHanSansSC-Bold.otf]{SourceHanSerifSC-Regular.otf}
\end{lstlisting}

这样写会覆盖模板的默认设置；如果只想调整某一个字形，请使用 \lstinline{fontspec} 提供的 \lstinline{ItalicFont}、\lstinline{BoldFont} 等选项，参见 \lstinline{texdoc fontspec}。

%% 第四章：长篇书稿的工程化组织——目录约定、\input 与 \include、按需加载宏包
\chapter{章节文件的组织与调用}
\label{chap:structure}

长篇书稿通常不止一个文件。把每一章放进 \lstinline{chapters/} 目录、由主文件统一调用，
有以下好处：单个文件短小，便于编辑与版本管理；可以只编译其中一部分；
也便于多人分工协作。

\section{目录与命名约定}

推荐的项目结构如下：

\begin{lstlisting}
BoyogalaElegantBook/
├── boyogalaelegantbook.cls   % 类文件
├── boyogalaelegantbook-cn.tex % 主文件：导言区 + 章节调用
├── chapters/                  % 各章内容
│   ├── chap-algorithm.tex     % 例：算法流程
│   ├── chap-graph.tex         % 例：图论
│   └── ...
├── fonts/                     % 字体文件
└── figure/, image/            % 插图
\end{lstlisting}

每个章节文件以 \lstinline{\\chapter} 开头，文件名建议使用统一的前缀
（如 \lstinline{chap-}）或按顺序编号（如 \lstinline{01-}、\lstinline{02-}），
以便排序与查阅。

\section{用 \texttt{\textbackslash input} 调用}

最简单的做法是在主文件中用 \lstinline{\\input} 引入章节文件：

\begin{lstlisting}
\begin{document}

\maketitle
\frontmatter
\tableofcontents
\mainmatter

% 依次调用各章；文件路径相对于主文件所在目录
% -----------------------------------------------------------------------------
% chapters/chap-algorithm.tex —— 算法流程案例：线性规划的单纯形法
%
% 本文件演示“把一章单独放进 chapters/ 目录、由主文件 \input 调用”的写法。
% 文件本身以 \chapter 开头，因此被调用后即成为书中的一章。
%
% 注意：本例只需要 booktabs 表格环境，不需要额外的算法宏包；
%       若要使用 algorithm2e、algorithmicx 等算法环境，请在主文件导言区加载，
%       而不是在类文件中加载。
% -----------------------------------------------------------------------------
\chapter{算法流程案例：线性规划的单纯形法}

本章以线性规划为例，说明如何把一个算法的执行流程写进书稿。算法流程的书写没有强制格式，
本模板提供的基本工具已经够用：\textbf{有序列表}用于描述步骤，\textbf{表格}用于展示中间结果，
\textbf{定理类环境}用于给出结论与证明，\lstinline{note} 环境用于补充说明。

\section{问题描述}

考虑如下线性规划问题：

\begin{equation}
  \begin{aligned}
    \max \quad & z = 3x_1 + 5x_2 \\
    \mathrm{s.t.}\quad & x_1 \leq 4, \\
                        & 2x_1 + 3x_2 \leq 18, \\
                        & x_1, x_2 \geq 0.
  \end{aligned}
  \label{eq:lp}
\end{equation}

目标是求使 $z$ 最大的可行解。可行域是一个多边形，最优解必然出现在某个顶点上，
但当约束较多时逐一枚举顶点并不可行，这正是单纯形法的用武之地。

\section{算法流程}

单纯形法的完整流程如下，每一步都对应表格的一次更新：

\begin{enumerate}
  \item \textbf{化为标准型}：引入松弛变量 $x_3, x_4$，把不等式约束改写为等式
    $x_1 + x_3 = 4$、$2x_1 + 3x_2 + x_4 = 18$，并把目标函数写为
    $z - 3x_1 - 5x_2 = 0$。以 $x_3, x_4$ 为初始基，写出初始单纯形表。
  \item \textbf{计算检验数}：观察目标函数行（$z$ 行）中各非基变量的系数。
    只要存在负系数，就说明当前解不是最优解。
  \item \textbf{选择入基变量}：取检验数最小（最负）的非基变量所在列。
    本例中 $x_2$ 的检验数为 $-5$，故 $x_2$ 入基。
  \item \textbf{计算比值并选择基变量}：对入基列中系数为正的各行计算比值
    $b_i / a_{ij}$，取最小者；该行对应的基变量离基。这一步保证主元为正，
    从而维持可行性。若某行该列系数不大于零，则不参与比值。
  \item \textbf{行变换（高斯消元）}：以主元为支点，把主元化为 $1$，再把同列其余元素
    化为 $0$，得到新的单纯形表。
  \item \textbf{判断是否最优}：若 $z$ 行不再有负系数，则当前基本可行解即为最优解；
    否则回到第 2 步继续迭代。
  \item \textbf{输出结果}：读出基变量取值与非基变量取值，代回目标函数得到最优值。
\end{enumerate}

\begin{note}
单纯形法的每一步迭代都保持可行性（比值检验的作用），但不能保证有限步终止——
退化情形下可能出现循环。因此实际应用中通常配合 \textbf{ Bland 规则} 等防循环策略。
\end{note}

\section{初始单纯形表与迭代}

初始时 $x_3, x_4$ 为基变量，非基变量为 $x_1, x_2$：

\begin{table}[htbp]
  \centering
  \caption{初始单纯形表}
  \label{tab:simplex-0}
  \begin{tabular}{c|ccccc}
    \toprule
    & $x_1$ & $x_2$ & $x_3$ & $x_4$ & $b$ \\
    \midrule
    $x_3$ & 1 & 0 & 1 & 0 & 4 \\
    $x_4$ & 2 & 3 & 0 & 1 & 18 \\
    \midrule
    $z$  & $-3$ & $-5$ & 0 & 0 & 0 \\
    \bottomrule
  \end{tabular}
\end{table}

$z$ 行中 $-5$ 是最小的负数，故 $x_2$ 入基。比值检验给出
$\frac{4}{0} = +\infty$（不参与）与 $\frac{18}{3} = 6$，故 $x_4$ 离基，
主元为第二行第三列的 $3$。行变换后得到：

\begin{table}[htbp]
  \centering
  \caption{迭代一次后的单纯形表}
  \label{tab:simplex-1}
  \begin{tabular}{c|ccccc}
    \toprule
    & $x_1$ & $x_2$ & $x_3$ & $x_4$ & $b$ \\
    \midrule
    $x_3$ & 1 & 0 & 1 & 0 & 4 \\
    $x_2$ & $\frac{2}{3}$ & 1 & 0 & $\frac{1}{3}$ & 6 \\
    \midrule
    $z$  & $\frac{1}{3}$ & 0 & 0 & $\frac{5}{3}$ & 30 \\
    \bottomrule
  \end{tabular}
\end{table}

此时 $z$ 行已无非负项以外的负数，迭代结束。

\section{结论}

\begin{conclusion}
问题 \eqref{eq:lp} 的最优解为 $x_1 = 0$、$x_2 = 6$，最优值为 $z = 30$。
\end{conclusion}

\begin{exercise}
用单纯形法求解下述线性规划，并写出所用的单纯形表：
\begin{equation}
  \max\; z = 2x_1 + 3x_2 \quad
  \mathrm{s.t.}\; x_1 + x_2 \leq 5,\; x_1 + 2x_2 \leq 8,\; x_1, x_2 \geq 0.
\end{equation}
\end{exercise}

% -----------------------------------------------------------------------------
% chapters/chap-graph.tex —— 图论案例：图的表示、度与最短路
%
% 本文件演示在章节中按需使用较重的宏包：画图需要 TikZ。
% 模板的类文件出于“精简依赖”的考虑不加载 TikZ，需要时请在主文件的导言区
% 写 \usepackage{tikz}；这样只有真正需要画图的文档才会付出编译代价。
% -----------------------------------------------------------------------------
\chapter{图论案例：图的表示、度与最短路}

\section{图与记号}

\begin{definition}[图]
设 $V$ 是一个非空有限集合，$E$ 是 $V$ 的若干二元子集构成的集合。
称有序对 $(V, E)$ 为一个\textbf{图}，记作 $G = (V, E)$，其中 $V$ 是图的
\textbf{顶点集}，$E$ 是图的\textbf{边集}；$E$ 中的每个二元子集称为一条边。
\end{definition}

\begin{remark}
若边 $\{u, v\}$ 的两个端点可以相同，则称 $u$ 与 $v$ 之间存在\textbf{自环}；
若 $E$ 中不含自环且任意两个顶点至多只相邻于一条边，则称 $G$ 为\textbf{简单图}。
本文讨论的图均为简单无向图。
\end{remark}

\section{度、路径与连通性}

\begin{definition}[度]
顶点 $v$ 的\textbf{度} $d(v)$ 是与 $v$ 相邻的顶点数；记 $\deg(G) = \sum_{v \in V} d(v)$。
\end{definition}

\begin{theorem}[握手定理]
任意有限图 $G$ 满足 $\sum_{v \in V} d(v) = 2|E|$。
\end{theorem}

\begin{proof}
对每条边计数：它连接两个顶点，对两个顶点的度数各贡献 $1$，
故对所有顶点的度数求和时每条边被计入两次。
\end{proof}

\begin{definition}[路径与连通]
顶点序列 $v_0, v_1, \dots, v_k$ 称为从 $v_0$ 到 $v_k$ 的一条\textbf{长度}为 $k$ 的
\textbf{路径}，若 $v_{i-1}$ 与 $v_i$ 相邻且各顶点互不相同。
若 $v_0 = v_k$，则称其为一个\textbf{圈}。若 $G$ 中任意两个顶点之间都存在路径，
则称 $G$ 是\textbf{连通}的。
\end{definition}

\section{图的一个例子}

考虑图 $G = (V, E)$，其中 $V = \{v_1, v_2, v_3, v_4, v_5\}$，
边集为 $E = \{\{v_1,v_2\}, \{v_2,v_3\}, \{v_1,v_4\}, \{v_4,v_5\}, \{v_5,v_3\}, \{v_2,v_5\}\}$，
各条边的权值如图 \ref{fig:graph-weighted} 所示。

\begin{figure}[htbp]
  \centering
  % 本图需要 TikZ，类文件不加载，请见主文件导言区的 \usepackage{tikz}
  \begin{tikzpicture}[
    scale=1.3,
    vertex/.style={circle,draw=structurecolor,fill=structurecolor!12,
                   inner sep=2pt,minimum size=6mm},
    weight/.style={fill=white,inner sep=1pt,font=\scriptsize}]
    \node[vertex] (v1) at (0,0) {$v_1$};
    \node[vertex] (v2) at (2.4,0) {$v_2$};
    \node[vertex] (v3) at (4.8,0) {$v_3$};
    \node[vertex] (v4) at (1.1,1.9) {$v_4$};
    \node[vertex] (v5) at (3.7,1.9) {$v_5$};
    \draw (v1) -- node[weight,above] {$1$} (v2);
    \draw (v2) -- node[weight,above] {$3$} (v3);
    \draw (v1) -- node[weight,left] {$2$} (v4);
    \draw (v4) -- node[weight,above] {$4$} (v5);
    \draw (v5) -- node[weight,right] {$1$} (v3);
    \draw (v2) -- node[weight,below right] {$5$} (v5);
  \end{tikzpicture}
  \caption{一个带权无向图}
  \label{fig:graph-weighted}
\end{figure}

\begin{exercise}
写出图 \ref{fig:graph-weighted} 中各顶点的度，并判断该图是否连通。
\end{exercise}

\section{图的表示}

计算机中存储图有两种常见方式：一是\textbf{邻接矩阵}，二是\textbf{邻接表}。
邻接矩阵用 $a_{ij} = 1$ 表示 $v_i$ 与 $v_j$ 相邻，用 $0$ 表示不相邻；
对于带权图，则把 $1$ 换成边权，不相邻的位置记为 $\infty$。

\begin{table}[htbp]
  \centering
  \caption{图 \ref{fig:graph-weighted} 的邻接矩阵}
  \label{tab:adjacency}
  \begin{tabular}{c|ccccc}
    \toprule
    & $v_1$ & $v_2$ & $v_3$ & $v_4$ & $v_5$ \\
    \midrule
    $v_1$ & 0 & 1 & 0 & 1 & 0 \\
    $v_2$ & 1 & 0 & 1 & 0 & 1 \\
    $v_3$ & 0 & 1 & 0 & 0 & 1 \\
    $v_4$ & 1 & 0 & 0 & 0 & 1 \\
    $v_5$ & 0 & 1 & 1 & 1 & 0 \\
    \bottomrule
  \end{tabular}
\end{table}

邻接矩阵需要 $O(|V|^2)$ 的存储空间，稠密图上运算方便；邻接表只需 $O(|V| + |E|)$ 的空间，
在稀疏图上更高效。

\section{最短路与遍历算法}

\begin{definition}[最短路]
在带权图中，路径 $P$ 的\textbf{长度}是其上各条边权之和。
顶点 $s$ 到 $t$ 的\textbf{最短路}是所有 $s$ 到 $t$ 的路径中长度最小者。
\end{definition}

\begin{custom}{算法：广度优先搜索}
从起点 $s$ 出发，逐层扩展：先访问所有与 $s$ 相邻的顶点，再访问与它们相邻的顶点，
如此进行，直到所有可达顶点都被访问。记录每个顶点首次被访问时所经过的边，
即得到一棵\textbf{广度优先搜索树}。
\end{custom}

\begin{note}
若所有边的权值相同（或者都取 $1$），广度优先搜索得到的就是最短路；
若边权各不相同，则需要 Dijkstra 算法等加权算法。
\end{note}

\begin{exercise}
在图 \ref{fig:graph-weighted} 上执行一次广度优先搜索，写出搜索树，
并求 $v_1$ 到 $v_3$ 的最短路长度。
\end{exercise}


\backmatter
\printbibliography[heading=bibintoc, title=\ebibname]

\end{document}
\end{lstlisting}

\lstinline{\\input} 的特点：

\begin{itemize}
  \item 只做文件展开，章节文件与主文件处于同一层级，\textbf{不能单独编译}；
  \item 章节文件中的 \lstinline{\\label} 与 \lstinline{\\ref} 可以跨章节互相引用；
  \item 一次编译即可完成，编译速度较快。
\end{itemize}

\section{用 \texttt{\textbackslash include} 调用}

如果希望每章可以单独编译（例如用 \lstinline{texdoc} 生成单章文档），
改用 \lstinline{\\include}：

\begin{lstlisting}
% -----------------------------------------------------------------------------
% chapters/chap-algorithm.tex —— 算法流程案例：线性规划的单纯形法
%
% 本文件演示“把一章单独放进 chapters/ 目录、由主文件 \input 调用”的写法。
% 文件本身以 \chapter 开头，因此被调用后即成为书中的一章。
%
% 注意：本例只需要 booktabs 表格环境，不需要额外的算法宏包；
%       若要使用 algorithm2e、algorithmicx 等算法环境，请在主文件导言区加载，
%       而不是在类文件中加载。
% -----------------------------------------------------------------------------
\chapter{算法流程案例：线性规划的单纯形法}

本章以线性规划为例，说明如何把一个算法的执行流程写进书稿。算法流程的书写没有强制格式，
本模板提供的基本工具已经够用：\textbf{有序列表}用于描述步骤，\textbf{表格}用于展示中间结果，
\textbf{定理类环境}用于给出结论与证明，\lstinline{note} 环境用于补充说明。

\section{问题描述}

考虑如下线性规划问题：

\begin{equation}
  \begin{aligned}
    \max \quad & z = 3x_1 + 5x_2 \\
    \mathrm{s.t.}\quad & x_1 \leq 4, \\
                        & 2x_1 + 3x_2 \leq 18, \\
                        & x_1, x_2 \geq 0.
  \end{aligned}
  \label{eq:lp}
\end{equation}

目标是求使 $z$ 最大的可行解。可行域是一个多边形，最优解必然出现在某个顶点上，
但当约束较多时逐一枚举顶点并不可行，这正是单纯形法的用武之地。

\section{算法流程}

单纯形法的完整流程如下，每一步都对应表格的一次更新：

\begin{enumerate}
  \item \textbf{化为标准型}：引入松弛变量 $x_3, x_4$，把不等式约束改写为等式
    $x_1 + x_3 = 4$、$2x_1 + 3x_2 + x_4 = 18$，并把目标函数写为
    $z - 3x_1 - 5x_2 = 0$。以 $x_3, x_4$ 为初始基，写出初始单纯形表。
  \item \textbf{计算检验数}：观察目标函数行（$z$ 行）中各非基变量的系数。
    只要存在负系数，就说明当前解不是最优解。
  \item \textbf{选择入基变量}：取检验数最小（最负）的非基变量所在列。
    本例中 $x_2$ 的检验数为 $-5$，故 $x_2$ 入基。
  \item \textbf{计算比值并选择基变量}：对入基列中系数为正的各行计算比值
    $b_i / a_{ij}$，取最小者；该行对应的基变量离基。这一步保证主元为正，
    从而维持可行性。若某行该列系数不大于零，则不参与比值。
  \item \textbf{行变换（高斯消元）}：以主元为支点，把主元化为 $1$，再把同列其余元素
    化为 $0$，得到新的单纯形表。
  \item \textbf{判断是否最优}：若 $z$ 行不再有负系数，则当前基本可行解即为最优解；
    否则回到第 2 步继续迭代。
  \item \textbf{输出结果}：读出基变量取值与非基变量取值，代回目标函数得到最优值。
\end{enumerate}

\begin{note}
单纯形法的每一步迭代都保持可行性（比值检验的作用），但不能保证有限步终止——
退化情形下可能出现循环。因此实际应用中通常配合 \textbf{ Bland 规则} 等防循环策略。
\end{note}

\section{初始单纯形表与迭代}

初始时 $x_3, x_4$ 为基变量，非基变量为 $x_1, x_2$：

\begin{table}[htbp]
  \centering
  \caption{初始单纯形表}
  \label{tab:simplex-0}
  \begin{tabular}{c|ccccc}
    \toprule
    & $x_1$ & $x_2$ & $x_3$ & $x_4$ & $b$ \\
    \midrule
    $x_3$ & 1 & 0 & 1 & 0 & 4 \\
    $x_4$ & 2 & 3 & 0 & 1 & 18 \\
    \midrule
    $z$  & $-3$ & $-5$ & 0 & 0 & 0 \\
    \bottomrule
  \end{tabular}
\end{table}

$z$ 行中 $-5$ 是最小的负数，故 $x_2$ 入基。比值检验给出
$\frac{4}{0} = +\infty$（不参与）与 $\frac{18}{3} = 6$，故 $x_4$ 离基，
主元为第二行第三列的 $3$。行变换后得到：

\begin{table}[htbp]
  \centering
  \caption{迭代一次后的单纯形表}
  \label{tab:simplex-1}
  \begin{tabular}{c|ccccc}
    \toprule
    & $x_1$ & $x_2$ & $x_3$ & $x_4$ & $b$ \\
    \midrule
    $x_3$ & 1 & 0 & 1 & 0 & 4 \\
    $x_2$ & $\frac{2}{3}$ & 1 & 0 & $\frac{1}{3}$ & 6 \\
    \midrule
    $z$  & $\frac{1}{3}$ & 0 & 0 & $\frac{5}{3}$ & 30 \\
    \bottomrule
  \end{tabular}
\end{table}

此时 $z$ 行已无非负项以外的负数，迭代结束。

\section{结论}

\begin{conclusion}
问题 \eqref{eq:lp} 的最优解为 $x_1 = 0$、$x_2 = 6$，最优值为 $z = 30$。
\end{conclusion}

\begin{exercise}
用单纯形法求解下述线性规划，并写出所用的单纯形表：
\begin{equation}
  \max\; z = 2x_1 + 3x_2 \quad
  \mathrm{s.t.}\; x_1 + x_2 \leq 5,\; x_1 + 2x_2 \leq 8,\; x_1, x_2 \geq 0.
\end{equation}
\end{exercise}

% -----------------------------------------------------------------------------
% chapters/chap-graph.tex —— 图论案例：图的表示、度与最短路
%
% 本文件演示在章节中按需使用较重的宏包：画图需要 TikZ。
% 模板的类文件出于“精简依赖”的考虑不加载 TikZ，需要时请在主文件的导言区
% 写 \usepackage{tikz}；这样只有真正需要画图的文档才会付出编译代价。
% -----------------------------------------------------------------------------
\chapter{图论案例：图的表示、度与最短路}

\section{图与记号}

\begin{definition}[图]
设 $V$ 是一个非空有限集合，$E$ 是 $V$ 的若干二元子集构成的集合。
称有序对 $(V, E)$ 为一个\textbf{图}，记作 $G = (V, E)$，其中 $V$ 是图的
\textbf{顶点集}，$E$ 是图的\textbf{边集}；$E$ 中的每个二元子集称为一条边。
\end{definition}

\begin{remark}
若边 $\{u, v\}$ 的两个端点可以相同，则称 $u$ 与 $v$ 之间存在\textbf{自环}；
若 $E$ 中不含自环且任意两个顶点至多只相邻于一条边，则称 $G$ 为\textbf{简单图}。
本文讨论的图均为简单无向图。
\end{remark}

\section{度、路径与连通性}

\begin{definition}[度]
顶点 $v$ 的\textbf{度} $d(v)$ 是与 $v$ 相邻的顶点数；记 $\deg(G) = \sum_{v \in V} d(v)$。
\end{definition}

\begin{theorem}[握手定理]
任意有限图 $G$ 满足 $\sum_{v \in V} d(v) = 2|E|$。
\end{theorem}

\begin{proof}
对每条边计数：它连接两个顶点，对两个顶点的度数各贡献 $1$，
故对所有顶点的度数求和时每条边被计入两次。
\end{proof}

\begin{definition}[路径与连通]
顶点序列 $v_0, v_1, \dots, v_k$ 称为从 $v_0$ 到 $v_k$ 的一条\textbf{长度}为 $k$ 的
\textbf{路径}，若 $v_{i-1}$ 与 $v_i$ 相邻且各顶点互不相同。
若 $v_0 = v_k$，则称其为一个\textbf{圈}。若 $G$ 中任意两个顶点之间都存在路径，
则称 $G$ 是\textbf{连通}的。
\end{definition}

\section{图的一个例子}

考虑图 $G = (V, E)$，其中 $V = \{v_1, v_2, v_3, v_4, v_5\}$，
边集为 $E = \{\{v_1,v_2\}, \{v_2,v_3\}, \{v_1,v_4\}, \{v_4,v_5\}, \{v_5,v_3\}, \{v_2,v_5\}\}$，
各条边的权值如图 \ref{fig:graph-weighted} 所示。

\begin{figure}[htbp]
  \centering
  % 本图需要 TikZ，类文件不加载，请见主文件导言区的 \usepackage{tikz}
  \begin{tikzpicture}[
    scale=1.3,
    vertex/.style={circle,draw=structurecolor,fill=structurecolor!12,
                   inner sep=2pt,minimum size=6mm},
    weight/.style={fill=white,inner sep=1pt,font=\scriptsize}]
    \node[vertex] (v1) at (0,0) {$v_1$};
    \node[vertex] (v2) at (2.4,0) {$v_2$};
    \node[vertex] (v3) at (4.8,0) {$v_3$};
    \node[vertex] (v4) at (1.1,1.9) {$v_4$};
    \node[vertex] (v5) at (3.7,1.9) {$v_5$};
    \draw (v1) -- node[weight,above] {$1$} (v2);
    \draw (v2) -- node[weight,above] {$3$} (v3);
    \draw (v1) -- node[weight,left] {$2$} (v4);
    \draw (v4) -- node[weight,above] {$4$} (v5);
    \draw (v5) -- node[weight,right] {$1$} (v3);
    \draw (v2) -- node[weight,below right] {$5$} (v5);
  \end{tikzpicture}
  \caption{一个带权无向图}
  \label{fig:graph-weighted}
\end{figure}

\begin{exercise}
写出图 \ref{fig:graph-weighted} 中各顶点的度，并判断该图是否连通。
\end{exercise}

\section{图的表示}

计算机中存储图有两种常见方式：一是\textbf{邻接矩阵}，二是\textbf{邻接表}。
邻接矩阵用 $a_{ij} = 1$ 表示 $v_i$ 与 $v_j$ 相邻，用 $0$ 表示不相邻；
对于带权图，则把 $1$ 换成边权，不相邻的位置记为 $\infty$。

\begin{table}[htbp]
  \centering
  \caption{图 \ref{fig:graph-weighted} 的邻接矩阵}
  \label{tab:adjacency}
  \begin{tabular}{c|ccccc}
    \toprule
    & $v_1$ & $v_2$ & $v_3$ & $v_4$ & $v_5$ \\
    \midrule
    $v_1$ & 0 & 1 & 0 & 1 & 0 \\
    $v_2$ & 1 & 0 & 1 & 0 & 1 \\
    $v_3$ & 0 & 1 & 0 & 0 & 1 \\
    $v_4$ & 1 & 0 & 0 & 0 & 1 \\
    $v_5$ & 0 & 1 & 1 & 1 & 0 \\
    \bottomrule
  \end{tabular}
\end{table}

邻接矩阵需要 $O(|V|^2)$ 的存储空间，稠密图上运算方便；邻接表只需 $O(|V| + |E|)$ 的空间，
在稀疏图上更高效。

\section{最短路与遍历算法}

\begin{definition}[最短路]
在带权图中，路径 $P$ 的\textbf{长度}是其上各条边权之和。
顶点 $s$ 到 $t$ 的\textbf{最短路}是所有 $s$ 到 $t$ 的路径中长度最小者。
\end{definition}

\begin{custom}{算法：广度优先搜索}
从起点 $s$ 出发，逐层扩展：先访问所有与 $s$ 相邻的顶点，再访问与它们相邻的顶点，
如此进行，直到所有可达顶点都被访问。记录每个顶点首次被访问时所经过的边，
即得到一棵\textbf{广度优先搜索树}。
\end{custom}

\begin{note}
若所有边的权值相同（或者都取 $1$），广度优先搜索得到的就是最短路；
若边权各不相同，则需要 Dijkstra 算法等加权算法。
\end{note}

\begin{exercise}
在图 \ref{fig:graph-weighted} 上执行一次广度优先搜索，写出搜索树，
并求 $v_1$ 到 $v_3$ 的最短路长度。
\end{exercise}

\end{lstlisting}

\lstinline{\\include} 会为每个文件建立独立的辅助文件与书签，代价是：

\begin{itemize}
  \item 编译速度略慢，每一章都要单独处理一遍交叉引用；
  \item 必须在导言区写 \lstinline{\\includeonly}，否则无法只编译部分章节；
  \item 被 \lstinline{\\includeonly} 排除的章节，引用其中的编号会显示为问号。
\end{itemize}

\begin{lstlisting}
% 只编译指定章节，其余章节跳过
\includeonly{chapters/chap-algorithm,chapters/chap-graph}
\end{lstlisting}

\section{在章节中按需加载宏包}

模板的类文件只加载排版所必需的宏包，尽量保持精简。像 TikZ、algorithm2e、longtable
这类较重或较专门的宏包，应当在\textbf{需要它们的主文件导言区}加载，
而不是让所有用户都付出编译代价：

\begin{lstlisting}
% 主文件导言区
\usepackage{tikz}          % 本书有画图章节，需要 TikZ
\end{lstlisting}

\begin{note}
在章节文件内部也可以写 \lstinline{\\usepackage}，但推荐放在导言区：
写在章节里会在文档正文中加载宏包，可能与 hyperref 等宏包产生选项冲突。
\end{note}

\section{示例章节}

下面两节分别给出两个完整的示例章节，它们就是上面 \lstinline{chapters/}
目录中的两个文件，由主文件用 \lstinline{\\input} 调用：

% -----------------------------------------------------------------------------
% chapters/chap-algorithm.tex —— 算法流程案例：线性规划的单纯形法
%
% 本文件演示“把一章单独放进 chapters/ 目录、由主文件 \input 调用”的写法。
% 文件本身以 \chapter 开头，因此被调用后即成为书中的一章。
%
% 注意：本例只需要 booktabs 表格环境，不需要额外的算法宏包；
%       若要使用 algorithm2e、algorithmicx 等算法环境，请在主文件导言区加载，
%       而不是在类文件中加载。
% -----------------------------------------------------------------------------
\chapter{算法流程案例：线性规划的单纯形法}

本章以线性规划为例，说明如何把一个算法的执行流程写进书稿。算法流程的书写没有强制格式，
本模板提供的基本工具已经够用：\textbf{有序列表}用于描述步骤，\textbf{表格}用于展示中间结果，
\textbf{定理类环境}用于给出结论与证明，\lstinline{note} 环境用于补充说明。

\section{问题描述}

考虑如下线性规划问题：

\begin{equation}
  \begin{aligned}
    \max \quad & z = 3x_1 + 5x_2 \\
    \mathrm{s.t.}\quad & x_1 \leq 4, \\
                        & 2x_1 + 3x_2 \leq 18, \\
                        & x_1, x_2 \geq 0.
  \end{aligned}
  \label{eq:lp}
\end{equation}

目标是求使 $z$ 最大的可行解。可行域是一个多边形，最优解必然出现在某个顶点上，
但当约束较多时逐一枚举顶点并不可行，这正是单纯形法的用武之地。

\section{算法流程}

单纯形法的完整流程如下，每一步都对应表格的一次更新：

\begin{enumerate}
  \item \textbf{化为标准型}：引入松弛变量 $x_3, x_4$，把不等式约束改写为等式
    $x_1 + x_3 = 4$、$2x_1 + 3x_2 + x_4 = 18$，并把目标函数写为
    $z - 3x_1 - 5x_2 = 0$。以 $x_3, x_4$ 为初始基，写出初始单纯形表。
  \item \textbf{计算检验数}：观察目标函数行（$z$ 行）中各非基变量的系数。
    只要存在负系数，就说明当前解不是最优解。
  \item \textbf{选择入基变量}：取检验数最小（最负）的非基变量所在列。
    本例中 $x_2$ 的检验数为 $-5$，故 $x_2$ 入基。
  \item \textbf{计算比值并选择基变量}：对入基列中系数为正的各行计算比值
    $b_i / a_{ij}$，取最小者；该行对应的基变量离基。这一步保证主元为正，
    从而维持可行性。若某行该列系数不大于零，则不参与比值。
  \item \textbf{行变换（高斯消元）}：以主元为支点，把主元化为 $1$，再把同列其余元素
    化为 $0$，得到新的单纯形表。
  \item \textbf{判断是否最优}：若 $z$ 行不再有负系数，则当前基本可行解即为最优解；
    否则回到第 2 步继续迭代。
  \item \textbf{输出结果}：读出基变量取值与非基变量取值，代回目标函数得到最优值。
\end{enumerate}

\begin{note}
单纯形法的每一步迭代都保持可行性（比值检验的作用），但不能保证有限步终止——
退化情形下可能出现循环。因此实际应用中通常配合 \textbf{ Bland 规则} 等防循环策略。
\end{note}

\section{初始单纯形表与迭代}

初始时 $x_3, x_4$ 为基变量，非基变量为 $x_1, x_2$：

\begin{table}[htbp]
  \centering
  \caption{初始单纯形表}
  \label{tab:simplex-0}
  \begin{tabular}{c|ccccc}
    \toprule
    & $x_1$ & $x_2$ & $x_3$ & $x_4$ & $b$ \\
    \midrule
    $x_3$ & 1 & 0 & 1 & 0 & 4 \\
    $x_4$ & 2 & 3 & 0 & 1 & 18 \\
    \midrule
    $z$  & $-3$ & $-5$ & 0 & 0 & 0 \\
    \bottomrule
  \end{tabular}
\end{table}

$z$ 行中 $-5$ 是最小的负数，故 $x_2$ 入基。比值检验给出
$\frac{4}{0} = +\infty$（不参与）与 $\frac{18}{3} = 6$，故 $x_4$ 离基，
主元为第二行第三列的 $3$。行变换后得到：

\begin{table}[htbp]
  \centering
  \caption{迭代一次后的单纯形表}
  \label{tab:simplex-1}
  \begin{tabular}{c|ccccc}
    \toprule
    & $x_1$ & $x_2$ & $x_3$ & $x_4$ & $b$ \\
    \midrule
    $x_3$ & 1 & 0 & 1 & 0 & 4 \\
    $x_2$ & $\frac{2}{3}$ & 1 & 0 & $\frac{1}{3}$ & 6 \\
    \midrule
    $z$  & $\frac{1}{3}$ & 0 & 0 & $\frac{5}{3}$ & 30 \\
    \bottomrule
  \end{tabular}
\end{table}

此时 $z$ 行已无非负项以外的负数，迭代结束。

\section{结论}

\begin{conclusion}
问题 \eqref{eq:lp} 的最优解为 $x_1 = 0$、$x_2 = 6$，最优值为 $z = 30$。
\end{conclusion}

\begin{exercise}
用单纯形法求解下述线性规划，并写出所用的单纯形表：
\begin{equation}
  \max\; z = 2x_1 + 3x_2 \quad
  \mathrm{s.t.}\; x_1 + x_2 \leq 5,\; x_1 + 2x_2 \leq 8,\; x_1, x_2 \geq 0.
\end{equation}
\end{exercise}


% -----------------------------------------------------------------------------
% chapters/chap-graph.tex —— 图论案例：图的表示、度与最短路
%
% 本文件演示在章节中按需使用较重的宏包：画图需要 TikZ。
% 模板的类文件出于“精简依赖”的考虑不加载 TikZ，需要时请在主文件的导言区
% 写 \usepackage{tikz}；这样只有真正需要画图的文档才会付出编译代价。
% -----------------------------------------------------------------------------
\chapter{图论案例：图的表示、度与最短路}

\section{图与记号}

\begin{definition}[图]
设 $V$ 是一个非空有限集合，$E$ 是 $V$ 的若干二元子集构成的集合。
称有序对 $(V, E)$ 为一个\textbf{图}，记作 $G = (V, E)$，其中 $V$ 是图的
\textbf{顶点集}，$E$ 是图的\textbf{边集}；$E$ 中的每个二元子集称为一条边。
\end{definition}

\begin{remark}
若边 $\{u, v\}$ 的两个端点可以相同，则称 $u$ 与 $v$ 之间存在\textbf{自环}；
若 $E$ 中不含自环且任意两个顶点至多只相邻于一条边，则称 $G$ 为\textbf{简单图}。
本文讨论的图均为简单无向图。
\end{remark}

\section{度、路径与连通性}

\begin{definition}[度]
顶点 $v$ 的\textbf{度} $d(v)$ 是与 $v$ 相邻的顶点数；记 $\deg(G) = \sum_{v \in V} d(v)$。
\end{definition}

\begin{theorem}[握手定理]
任意有限图 $G$ 满足 $\sum_{v \in V} d(v) = 2|E|$。
\end{theorem}

\begin{proof}
对每条边计数：它连接两个顶点，对两个顶点的度数各贡献 $1$，
故对所有顶点的度数求和时每条边被计入两次。
\end{proof}

\begin{definition}[路径与连通]
顶点序列 $v_0, v_1, \dots, v_k$ 称为从 $v_0$ 到 $v_k$ 的一条\textbf{长度}为 $k$ 的
\textbf{路径}，若 $v_{i-1}$ 与 $v_i$ 相邻且各顶点互不相同。
若 $v_0 = v_k$，则称其为一个\textbf{圈}。若 $G$ 中任意两个顶点之间都存在路径，
则称 $G$ 是\textbf{连通}的。
\end{definition}

\section{图的一个例子}

考虑图 $G = (V, E)$，其中 $V = \{v_1, v_2, v_3, v_4, v_5\}$，
边集为 $E = \{\{v_1,v_2\}, \{v_2,v_3\}, \{v_1,v_4\}, \{v_4,v_5\}, \{v_5,v_3\}, \{v_2,v_5\}\}$，
各条边的权值如图 \ref{fig:graph-weighted} 所示。

\begin{figure}[htbp]
  \centering
  % 本图需要 TikZ，类文件不加载，请见主文件导言区的 \usepackage{tikz}
  \begin{tikzpicture}[
    scale=1.3,
    vertex/.style={circle,draw=structurecolor,fill=structurecolor!12,
                   inner sep=2pt,minimum size=6mm},
    weight/.style={fill=white,inner sep=1pt,font=\scriptsize}]
    \node[vertex] (v1) at (0,0) {$v_1$};
    \node[vertex] (v2) at (2.4,0) {$v_2$};
    \node[vertex] (v3) at (4.8,0) {$v_3$};
    \node[vertex] (v4) at (1.1,1.9) {$v_4$};
    \node[vertex] (v5) at (3.7,1.9) {$v_5$};
    \draw (v1) -- node[weight,above] {$1$} (v2);
    \draw (v2) -- node[weight,above] {$3$} (v3);
    \draw (v1) -- node[weight,left] {$2$} (v4);
    \draw (v4) -- node[weight,above] {$4$} (v5);
    \draw (v5) -- node[weight,right] {$1$} (v3);
    \draw (v2) -- node[weight,below right] {$5$} (v5);
  \end{tikzpicture}
  \caption{一个带权无向图}
  \label{fig:graph-weighted}
\end{figure}

\begin{exercise}
写出图 \ref{fig:graph-weighted} 中各顶点的度，并判断该图是否连通。
\end{exercise}

\section{图的表示}

计算机中存储图有两种常见方式：一是\textbf{邻接矩阵}，二是\textbf{邻接表}。
邻接矩阵用 $a_{ij} = 1$ 表示 $v_i$ 与 $v_j$ 相邻，用 $0$ 表示不相邻；
对于带权图，则把 $1$ 换成边权，不相邻的位置记为 $\infty$。

\begin{table}[htbp]
  \centering
  \caption{图 \ref{fig:graph-weighted} 的邻接矩阵}
  \label{tab:adjacency}
  \begin{tabular}{c|ccccc}
    \toprule
    & $v_1$ & $v_2$ & $v_3$ & $v_4$ & $v_5$ \\
    \midrule
    $v_1$ & 0 & 1 & 0 & 1 & 0 \\
    $v_2$ & 1 & 0 & 1 & 0 & 1 \\
    $v_3$ & 0 & 1 & 0 & 0 & 1 \\
    $v_4$ & 1 & 0 & 0 & 0 & 1 \\
    $v_5$ & 0 & 1 & 1 & 1 & 0 \\
    \bottomrule
  \end{tabular}
\end{table}

邻接矩阵需要 $O(|V|^2)$ 的存储空间，稠密图上运算方便；邻接表只需 $O(|V| + |E|)$ 的空间，
在稀疏图上更高效。

\section{最短路与遍历算法}

\begin{definition}[最短路]
在带权图中，路径 $P$ 的\textbf{长度}是其上各条边权之和。
顶点 $s$ 到 $t$ 的\textbf{最短路}是所有 $s$ 到 $t$ 的路径中长度最小者。
\end{definition}

\begin{custom}{算法：广度优先搜索}
从起点 $s$ 出发，逐层扩展：先访问所有与 $s$ 相邻的顶点，再访问与它们相邻的顶点，
如此进行，直到所有可达顶点都被访问。记录每个顶点首次被访问时所经过的边，
即得到一棵\textbf{广度优先搜索树}。
\end{custom}

\begin{note}
若所有边的权值相同（或者都取 $1$），广度优先搜索得到的就是最短路；
若边权各不相同，则需要 Dijkstra 算法等加权算法。
\end{note}

\begin{exercise}
在图 \ref{fig:graph-weighted} 上执行一次广度优先搜索，写出搜索树，
并求 $v_1$ 到 $v_3$ 的最短路长度。
\end{exercise}


%% 第五章：综合演示——定理环境、交叉引用、图表与数学公式
\chapter{BoyogalaElegantBook 写作示例}

\begin{introduction}
  \item 积分定义~\ref{def:int}
  \item Fubini 定理~\ref{thm:fubi}
  \item 最优性原理~\ref{pro:max}
  \item 柯西列性质~\ref{property:cauchy}
  \item 韦达定理
\end{introduction}

\section{Lebesgue 积分}
在前面各章做了必要的准备后，本章开始介绍新的积分。在 Lebesgue 测度理论的基础上建立了 Lebesgue 积分，其被积函数和积分域更一般，可以对有界函数和无界函数统一处理。正是由于 Lebesgue 积分的这些特点，使得 Lebesgue 积分比 Riemann 积分具有在更一般条件下的极限定理和累次积分交换积分顺序的定理，这使得 Lebesgue 积分不仅在理论上更完善，而且在计算上更灵活有效。

Lebesgue 积分有几种不同的定义方式。我们将采用逐步定义非负简单函数，非负可测函数和一般可测函数积分的方式。

由于现代数学的许多分支如概率论、泛函分析、调和分析等常常用到一般空间上的测度与积分理论，在本章最后一节将介绍一般的测度空间上的积分。

\subsection{积分的定义}

我们将通过三个步骤定义可测函数的积分。首先定义非负简单函数的积分。以下设 $E$ 是 $\mathcal{R}^n$ 中的可测集。

\begin{definition}[可积性] \label{def:int} 
设 $ f(x)=\sum\limits_{i=1}^{k} a_i \chi_{A_i}(x)$ 是 $E$ 上的\textbf{非负简单函数}，中文其中 $\{A_1,A_2,\ldots,A_k\}$ 是 $E$ 上的一个可测分割，$a_1,a_2,\ldots,a_k$ 是非负实数。定义 $f$ 在 $E$ 上的积分为 $\int_{a}^b f(x)$
\begin{equation}
   \label{inter}
   \int_{E} f dx = \sum_{i=1}^k a_i m(A_i) \pi \alpha\beta\sigma\gamma\nu\xi\epsilon\varepsilon. \oint_{a}^b\prod_{i=1}^n
\end{equation}
一般情况下 $0 \leq \int_{E} f dx \leq \infty$。若 $\int_{E} f dx < \infty$，则称 $f$ 在 $E$ 上可积。
\end{definition}

一个自然的问题是，Lebesgue 积分与我们所熟悉的 Riemann 积分有什么联系和区别？在 4.4 在我们将详细讨论 Riemann 积分与 Lebesgue 积分的关系。这里只看一个简单的例子。设 $D(x)$ 是区间 $[0,1]$ 上的 Dirichlet 函数。即 $D(x)=\chi_{Q_0}(x)$，其中 $Q_0$ 表示 $[0,1]$ 中的有理数的全体。根据非负简单函数积分的定义，$D(x)$ 在 $[0,1]$ 上的 Lebesgue 积分为
\begin{equation}
   \label{inter2}
   \int_0^1 D(x)dx = \int_0^1 \chi_{Q_0} (x) dx = m(Q_0) = 0
\end{equation}
即 $D(x)$ 在 $[0,1]$ 上是 Lebesgue 可积的并且积分值为零。但 $D(x)$ 在 $[0,1]$ 上不是 Riemann 可积的。


有界变差函数是与单调函数有密切联系的一类函数。有界变差函数可以表示为两个单调递增函数之差。与单调函数一样，有界变差函数几乎处处可导。与单调函数不同，有界变差函数类对线性运算是封闭的，它们构成一线空间。练习题 \ref{exer:43} 是一个性质的证明。

\begin{exercise}\label{exer:43}
设 $f \notin\in L(\mathcal{R}^1)$，$g$ 是 $\mathcal{R}^1$ 上的有界可测函数。证明函数
\begin{equation}
   \label{ex:1}
   I(t) = \int_{\mathcal{R}^1} f(x+t)g(x)dx \quad t \in \mathcal{R}^1
\end{equation}
是 $\mathcal{R}^1$ 上的连续函数。 
\end{exercise}

\begin{solution}
即 $D(x)$ 在 $[0,1]$ 上是 Lebesgue 可积的并且积分值为零。但 $D(x)$ 在 $[0,1]$ 上不是 Riemann 可积的。
\end{solution}

\begin{proof}
即 $D(x)$ 在 $[0,1]$ 上是 Lebesgue 可积的并且积分值为零。但 $D(x)$ 在 $[0,1]$ 上不是 Riemann 可积的。
\end{proof}

\begin{theorem}[Fubini 定理] \label{thm:fubi} 
（1）若 $f(x,y)$ 是 $\mathcal{R}^p\times\mathcal{R}^q$ 上的非负可测函数，则对几乎处处的 $x\in \mathcal{R}^p$，$f(x,y)$ 作为 $y$ 的函数是 $\mathcal{R}^q$ 上的非负可测函数，$g(x)=\int_{\mathcal{R}^q}f(x,y) dy$ 是 $\mathcal{R}^p$ 上的非负可测函数。并且
\begin{equation}
   \label{eq:461}
   \int_{\mathcal{R}^p\times\mathcal{R}^q} f(x,y) dxdy=\int_{\mathcal{R}^p}\left(\int_{\mathcal{R}^q}f(x,y)dy\right)dx.
\end{equation}

（2）若 $f(x,y)$ 是 $\mathcal{R}^p\times\mathcal{R}^q$ 上的可积函数，则对几乎处处的 $x\in\mathcal{R}^p$，$f(x,y)$ 作为 $y$ 的函数是 $\mathcal{R}^q$ 上的可积函数，并且 $g(x)=\int_{\mathcal{R}^q}f(x,y) dy$ 是 $\mathcal{R}^p$ 上的可积函数。而且~\eqref{eq:461} 成立。
\end{theorem}

\begin{note}
在本模板中，引理（lemma），推论（corollary）的样式和定理~\ref{thm:fubi} 的样式一致，包括颜色，仅仅只有计数器的设置不一样。
\end{note}

我们说一个实变或者复变量的实值或者复值函数是在区间上平方可积的，如果其绝对值的平方在该区间上的积分是有限的。所有在勒贝格积分意义下平方可积的可测函数构成一个希尔伯特空间，也就是所谓的 $L^2$ 空间，几乎处处相等的函数归为同一等价类。形式上，$L^2$ 是平方可积函数的空间和几乎处处为 0 的函数空间的商空间。

\begin{proposition}[最优性原理] \label{pro:max}
如果 $u^*$ 在 $[s,T]$ 上为最优解，则 $u^*$ 在 $[s, T]$ 任意子区间都是最优解，假设区间为 $[t_0, t_1]$ 的最优解为 $u^*$ ，则 $u(t_0)=u^{*}(t_0)$，即初始条件必须还是在 $u^*$ 上。
\end{proposition}

我们知道最小二乘法可以用来处理一组数据，可以从一组测定的数据中寻求变量之间的依赖关系，这种函数关系称为经验公式。本课题将介绍最小二乘法的精确定义及如何寻求点与点之间近似成线性关系时的经验公式。假定实验测得变量之间的 $n$ 个数据，则在平面上，可以得到 $n$ 个点，这种图形称为 “散点图”，从图中可以粗略看出这些点大致散落在某直线近旁, 我们认为其近似为一线性函数，下面介绍求解步骤。

\begin{figure}[htbp]
  \centering
  \includegraphics[width=0.6\textwidth]{scatter.jpg}
  \caption{散点图示例 $\hat{y}=a+bx$ \label{fig:scatter}}
\end{figure}

以最简单的一元线性模型来解释最小二乘法。什么是一元线性模型呢？监督学习中，如果预测的变量是离散的，我们称其为分类（如决策树，支持向量机等），如果预测的变量是连续的，我们称其为回归。回归分析中，如果只包括一个自变量和一个因变量，且二者的关系可用一条直线近似表示，这种回归分析称为一元线性回归分析。如果回归分析中包括两个或两个以上的自变量，且因变量和自变量之间是线性关系，则称为多元线性回归分析。对于二维空间线性是一条直线；对于三维空间线性是一个平面，对于多维空间线性是一个超平面。

\begin{property}\label{property:cauchy}
柯西列的性质
\begin{enumerate}
\item $\{x_k\}$ 是柯西列，则其子列 $\{x_k^i\}$ 也是柯西列。
\item $x_k\in \mathcal{R}^n$，$\rho(x,y)$ 是欧几里得空间，则柯西列收敛，$(\mathcal{R}^n,\rho)$ 空间是完备的。
\end{enumerate}
\end{property}

\begin{conclusion}
回归分析（regression analysis） 是确定两种或两种以上变量间相互依赖的定量关系的一种统计分析方法。运用十分广泛，回归分析按照涉及的变量的多少，分为一元回归和多元回归分析；按照因变量的多少，可分为简单回归分析和多重回归分析；按照自变量和因变量之间的关系类型，可分为线性回归分析和非线性回归分析。
\end{conclusion}

\begin{problemset}
\item 设 $A$ 为数域 $K$ 上的 $n$ 级矩阵。证明：如果 $K^n$ 中任意非零列向量都是 $A$ 的特征向量，则 $A$ 一定是数量矩阵。
\item 证明：不为零矩阵的幂零矩阵不能对角化。
\item 设 $A = (a_{ij})$ 是数域 $K$ 上的一个 $n$ 级上三角矩阵，证明：如果 $a_{11} = a_{22} = \cdots = a_{nn}$，并且至少有一个 $a_{kl} \not = 0 (k < l)$，则 $A$ 一定不能对角化。
\end{problemset}

%% 第六章：用户社区反馈的常见问题
\chapter{常见问题集}

我们根据用户社区反馈整理了下面一些常见的问题，用户在遇到问题时，应当首先查阅本手册和本部分的常见的问题。

\begin{enumerate}[itemsep=1.5ex]
  \item \question{有没有办法章节用“第一章，第一节，（一）”这种？}
    见前文介绍，可以使用 \lstinline{scheme=chinese} 设置。
  \item \question{大佬，我想把正文字体改为亮色，背景色改为黑灰色。}
    页面颜色可以使用 \lstinline{\pagecolor} 命令设置，文本命令可以参考\href{https://tex.stackexchange.com/questions/278544/xcolor-what-is-the-equivalent-of-default-text-color}{这里}进行设置。
  \item \question{\lstinline{! LaTeX Error: Unknown option 'scheme=plain' for package 'ctex'.}}
    你用的 \CTeX\ 套装吧？这个里面的 \lstinline{ctex} 宏包已经是已经是 10 年前的了，与本模板使用的 \lstinline{ctex} 宏集有很大区别。不建议 \CTeX\ 套装了，请卸载并安装 \TeX~Live 2025。
  \item \question{我该使用什么版本？}
    请务必使用\href{https://github.com/boyogala/BoyogalaElegantBook/releases}{最新正式发行版}，发行版间不定期可能会有更新（修复 bug 或者改进之类），如果你在使用过程中没有遇到问题，不需要每次更新\href{https://github.com/boyogala/BoyogalaElegantBook/archive/refs/heads/main.zip}{最新版}，但是在发行版更新之后，请尽可能使用最新版（发行版）！最新发行版可以在 GitHub 的 Releases 页面获取。
  \item \question{我该使用什么编辑器？}
    你可以使用 \TeX~Live 2025 自带的编辑器 \TeX works 或者使用 \TeX studio，\TeX works 的自动补全，你可以参考我们的总结 \href{https://github.com/EthanDeng/texworks-autocomplete}{\TeX works 自动补全}。推荐使用 \TeX~Live 2025 + \TeX studio。我（注：此为前维护者）自己用 VS Code 和 Sublime Text，相关的配置说明，请参考 \href{https://github.com/EthanDeng/vscode-latex}{\LaTeX\ 编译环境配置：Visual Studio Code 配置简介} 和 \href{https://github.com/EthanDeng/sublime-text-latex}{Sublime Text 搭建 \LaTeX\ 编写环境}。
  \item \question{您好，我们想用您的 BoyogalaElegantBook 模板写一本书。关于机器学习的教材，希望获得您的授权，谢谢您的宝贵时间。}
    模板的使用修改都是自由的，你们声明模板来源以及模板地址（GitHub 地址）即可，其他未尽事宜按照开源协议 LPPL-1.3c。做好之后，如果方便的话，可以给我们一个链接，我把你们的教材放在 Elegant\LaTeX\ 用户作品集里。
  \item \question{请问交叉引用是什么？}
    本群和本模板适合有一定 \LaTeX\ 基础的用户使用，新手请先学习 \LaTeX\ 的基础，理解各种概念，否则你将寸步难行。
  \item \question{代码高亮环境能用其他语言吗？}
    可以的，BoyogalaElegantBook 模板用的是 \lstinline{listings} 宏包，你可以在环境（\lstinline{lstlisting}）之后加上语言（比如 Python 使用 \lstinline{language=Python} 选项），全局语言修改请使用 \lstinline{lstset} 命令，更多信息请参考宏包文档。
  \item \question{群主，什么时候出 Beamer 的模板（主题），ElegantSlide 或者 ElegantBeamer？}
    由于 Beamer 中有一个很优秀的主题 \href{https://github.com/matze/mtheme}{Metropolis}。后续确定不会再出任何主题/模板，请大家根据需要修改已有主题。
\end{enumerate}

%% 第七章：修订历史，由 \datechange 与 change 环境生成
\chapter{版本更新历史}

根据用户的反馈，我们不断修正和完善模板。由于 3.00 之前版本与现在版本差异非常大，在此不列出 3.00 之前的更新内容。

\datechange{2026/10/22}{BoyogalaElegantBook 首个版本}
\begin{change}
  \item 项目更名为 BoyogalaElegantBook，类文件更名为 \lstinline{boyogalaelegantbook.cls}；
  \item 只保留一种字体配置：西文、中文与数学字体全部随项目提供于 \lstinline{fonts/} 目录，均为 SIL OFL 1.1；
  \item 仅支持 XeLaTeX，移除 \lstinline{math}、\lstinline{chinesefont} 等字体选项；
  \item 精简依赖：移除 \lstinline{tikz}、\lstinline{mwe}、\lstinline{lipsum}、\lstinline{makecell}、\lstinline{hologo} 等宏包，装饰符号改用数学字体字形；
  \item 修复 \lstinline{color=none} 时的未定义颜色错误；语言模式精简为简体中文、繁体中文与英文。
\end{change}

\datechange{2026/02/27}{发布 v4.6 版本}  \textcolor{blue}{\bfseries 恢复维护}

\begin{change}
  \item 确认新维护者啸行并恢复维护。
  \item 鉴于用户群体庞大，模板重新发布。
  \item 更新以支持 \TeX~Live 2025 及当前主流操作系统（Windows 11、Ubuntu 24.04 LTS、macOS）。
\end{change}

\datechange{2022/12/31}{版本 4.5} \textcolor{red}{\bfseries 停止维护！}

\datechange{2022/08/17}{版本 4.5 pre}
\begin{change}
  \item \textbf{重要改动}：提供了一个新的文档类选项 \lstinline|usesamecnt|，可以使全局的定理类环境使用同一个计数器。
  \item \textbf{重要改动}：修改了 \lstinline|\elegantnewtheorem| 命令，使其有第五个（可选）参数。
\end{change}

\datechange{2022/08/15}{版本 4.4 正式发布。}

\begin{change}
  \item \textbf{重要改动}：提供了一个定义定理类环境的命令 \lstinline|\elegantnewtheorem|；
  \item \textbf{重要改动}：为所有内置定理类环境提供了带星号的版本，带星号的定理类环境不会编号，修复 \href{https://github.com/ElegantLaTeX/ElegantBook/issues/167}{issue: \#167}；
  \item \textbf{重要改动}：在 \lstinline{scheme=chinese} 下将目录中的“第 1 章”修改为“第一章”；
  \item 将 TeX Gyre Termes 改为 TeX Gyre TermesX，使英文部分字形与 newtx 系列宏包更相近；
  \item 重写了内置定理类环境的实现方法，修复了一些 bug，由于修改部分较大，如果引入了新的 bug，请及时在 QQ 群或 \href{https://github.com/ElegantLaTeX}{Github} 上进行反馈；
  \item 删除 Gitee 仓库地址，恢复 GitHub 提交（pull requests）；
  \item 将参考文献命令添加到导言区，使编辑器能够对参考文献自动补全。
\end{change}

\datechange{2022/04/09}{版本 4.3 正式发布。}

\begin{change}
  \item 放弃 newtx 系列宏包的设置，改用 TeX Gyre Termes，并设置其他字体；
  \item 修改定理类环境内部字体设置，修复环境内部中文无法加粗问题；
  \item 增加定理类环境的计数器选项 \lstinline{thmcnt}，可选 \lstinline{chapter} 和 \lstinline{section}；
  \item 增加 \lstinline{bibend} 选项，可选 \lstinline{bibend=biber}（默认）和 \lstinline{bibend=bibtex}。
\end{change}



\datechange{2022/03/08}{版本 4.2 正式发布。}

\begin{change}
  \item 对于 newtx 系列宏包更新导致的字体 bug 的修复；
  \item 修缮目录格式，为了达到这个目的，重新改写 \lstinline{\chaptername} 的重定义语句；
  \item 增加日语 \lstinline{lang=jp} 设定。
  \item 这个版本为一个临时性版本，在 \TeX Live 2022 发布之后，将尽快发布 4.3 版本，由于对于中文的改动比较大，可能会出现预期之外的 bug，有问题可以在 QQ 群或者 Github 反馈。
\end{change}


\datechange{2021/05/02}{版本 4.1 正式发布。}

\begin{change}
  \item \textbf{重要改动}：由原先的 \textsf{BibTeX} 改为 biblatex 编译方式（后端为 \lstinline{biber}），请注意两者之间的差异；
  \item \textbf{重要改进}：修改对于定理写法兼容方式，提高数学公式代码的兼容性；
  \item 页面设置改动，默认页面更宽；方便书写和阅读；
  \item 支持目录文字以及页码跳转；
  \item 不再维护 \textsf{pdfLaTeX} 中文支持方式，请务必使用 \textsf{XeLaTeX} 编译中文文稿。
  \item 增加多个语言选项，法语 \lstinline{lang=fr}、荷兰语 \lstinline{lang=nl}、匈牙利语 \lstinline{lang=hu}、西班牙语 \lstinline{lang=es}、蒙古语 \lstinline{lang=mn} 等。
\end{change}


\datechange{2020/04/12}{版本 3.11 正式发布，\textcolor{red}{此版本为 3.x 最后版本。}}

\begin{change}
  \item \textbf{重要修正}：修复因为 \lstinline{gbt7714} 宏包更新导致的 \lstinline{natbib option clash} 错误；
  \item 由于 \lstinline{pgfornament} 宏包未被 \TeX~Live 2020 收录，因此删除 base 相关的内容；
  \item 修复部分环境的空格问题；
  \item 增加了意大利语言选项 \lstinline{lang=it}。
\end{change}


\datechange{2020/02/10}{版本 3.10 正式发布}

\begin{change}
  \item 增加数学字体选项 \lstinline{math}，可选项为 \lstinline{newtx} 和 \lstinline{cm}。\\
  \textbf{重要提示}：原先通过 \lstinline{newtxmath} 宏包设置的数学字体改为 \LaTeX\ 默认数学字体，如果需要保持原来的字体，需要显式声明数学字体（\lstinline{math=newtx}）；
  \item 新增中文字体选项 \lstinline{chinesefont}，可选项为 \lstinline{ctexfont}、\lstinline{founder} 和 \lstinline{nofont}。
  \item 将封面作者信息设置为可选，并且增加自定义信息命令 \lstinline{\bioinfo}；
  \item 在说明文档中增加版本历史，新增 \lstinline{\datechange} 命令和 \lstinline{change} 环境；
  \item 增加汉化章节选项 \lstinline{scheme}，可选项为汉化 \lstinline{chinese}；
  \item 由于 \lstinline{\lvert} 问题已经修复，重新调整 \lstinline{ctex} 宏包和 \lstinline{amsmath} 宏包位置。
  \item 修改页眉设置，去除了 \lstinline{\lastpage} 以避免 page anchor 问题，加入 \lstinline{\frontmatter}。
  \item 修改参考文献选项 \lstinline{cite}，可选项为数字 \lstinline{numbers}、 作者-年份 \lstinline{authoryear} 以及上标 \lstinline{super}。
  \item 新增参考文献样式选项 \lstinline{bibstyle}，并将英文模式下参考文献样式 \lstinline{apalike} 设置为默认值，中文仍然使用 \lstinline{gbt7714} 宏包设置。
\end{change}

\datechange{2019/08/18}{版本 3.09 正式发布}

\begin{change}
  \item \lstinline{\elegantpar} 存在 bug，删除 \lstinline{\elegantpar} 命令，建议用户改用 \lstinline{\marginnote} 和 \lstinline{\marginpar} 旁注命令。
  \item 积分操作符统一更改为 \lstinline{esint} 宏包设置；
  \item 新增目录选项 \lstinline{toc}，可选项为单栏 \lstinline{onecol} 和双栏 \lstinline{twocol}；
  \item 手动增加参考文献选项 \lstinline{cite}，可选项为上标形式 \lstinline{super}；
  \item 修正章节习题（\lstinline{problemset}）环境。
\end{change}

\datechange{2019/05/28}{版本 3.08 正式发布}

\begin{change}
  \item 修复 \lstinline{\part} 命令。
  \item 引入 Note 模板中的 \lstinline{pad} 选项 \lstinline{device=pad}。
  \item 数学字体加入 \lstinline{mtpro2} 可选项 \lstinline{math=mtpro2}，使用免费的 \lstinline{lite} 子集。
  \item 将参考文献默认显示方式 \lstinline{authoyear} 改为 \lstinline{numbers}。
  \item 引入旁注命令 \lstinline{\marginpar}（测试）。
  \item 新增章节摘要环境 \lstinline{introduction}。
  \item 新增章节习题环境 \lstinline{problemset}。
  \item 将 \lstinline{\equote} 重命名为 \lstinline{\extrainfo}。
  \item 完善说明文档，增加致谢部分。
\end{change}

\datechange{2019/04/15}{版本 3.07 正式发布}

\begin{change}
  \item 删除中英文自定义字体总设置。
  \item 新增颜色主题，并将原绿色默认主题设置为蓝色 \lstinline{color=blue}。
  \item 引入隐藏装饰图案选项 \lstinline{base}，可选项有显示 \lstinline{show} 和隐藏 \lstinline{hide}。
  \item 新增定理模式 \lstinline{mode}，可选项有简单模式 \lstinline{simple} 和炫彩模式 \lstinline{fancy}。
  \item 新增隐藏证明、答案等环境的选项 \lstinline{result=noanswer}。
\end{change}

\datechange{2019/02/25}{版本 3.06 正式发布}

\begin{change}
  \item 删除水印。
  \item 新封面，新装饰图案。
  \item 添加引言使用说明。
  \item 修复双面 \lstinline{twoside}。
  \item 美化列表环境。
  \item 增加 \lstinline{\subsubsection} 的设置。
  \item 将模板拆分成中英文语言模式。
  \item 使用 \lstinline{lstlisting} 添加代码高亮。
  \item 增加定理类环境使用说明。
\end{change}

\datechange{2019/01/22}{版本 3.05 正式发布}

\begin{change}
  \item 添加 \lstinline{xeCJK} 宏包中文支持方案。
  \item 修复模板之前对 Ti\textit{k}Z 单位的改动。
  \item 更新 logo 图。
\end{change}

\datechange{2019/01/15}{版本 3.04 正式发布}

\begin{change}
  \item 格式化模板代码。
  \item 增加 \lstinline{\equote} 命令。
  \item 修改 \lstinline{\date}。
\end{change}

\datechange{2019/01/08}{版本 3.03 正式发布}

\begin{change}
  \item 修复附录章节显示问题。
  \item 小幅优化封面代码。
\end{change}

\datechange{2018/12/31}{版本 3.02 正式发布}

\begin{change}
  \item 修复名字系列命令自定义格式时出现的空格问题，比如 \lstinline{\listfigurename}。
  \item 英文定理类名字改为中文名。
  \item 英文结构名改为中文。
\end{change}

\datechange{2018/12/16}{版本 3.01 正式发布}

\begin{change}
  \item 调整 \lstinline{ctex} 宏包。
  \item 说明文档增加更新内容。
\end{change}

\datechange{2018/12/06}{版本 3.00 正式发布}

\begin{change}
  \item 删除 \lstinline{mathpazo} 数学字体选项。
  \item 添加邮箱命令 \lstinline{\mailto}。
  \item 修改英文字体为 \lstinline{newtx} 系列，另外大型操作符号维持 cm 字体。
  \item 中文字体改用 \lstinline{ctex} 宏包自动设置。
  \item 删除 \lstinline{xeCJK} 字体设置，原因是不同系统字体不方便统一。
  \item 定理换用 \lstinline{tcolorbox} 宏包定义，并基本维持原有的定理样式，优化显示效果，支持跨页；定理类名字重命名，如 etheorem 改为 theorem 等等。
  \item 删去自定义的缩进命令 \lstinline{\Eindent}。
  \item 添加参考文献宏包 \lstinline{natbib}。
  \item 颜色名字重命名。
\end{change}

\nocite{*}

\printbibliography[heading=bibintoc, title=\ebibname]
\appendix

%% 附录：书中用到的基本数学工具
\chapter{基本数学工具}


本附录包括了计量经济学中用到的一些基本数学，我们扼要论述了求和算子的各种性质，研究了线性和某些非线性方程的性质，并复习了比例和百分数。我们还介绍了一些在应用计量经济学中常见的特殊函数，包括二次函数和自然对数，前 4 节只要求基本的代数技巧，第 5 节则对微分学进行了简要回顾；虽然要理解本书的大部分内容，微积分并非必需，但在一些章末附录和第 3 篇某些高深专题中，我们还是用到了微积分。

\section{求和算子与描述统计量}

\textbf{求和算子} 是用以表达多个数求和运算的一个缩略符号，它在统计学和计量经济学分析中扮演着重要作用。如果 $\{x_i: i=1, 2, \ldots, n\}$ 表示 $n$ 个数的一个序列，那么我们就把这 $n$ 个数的和写为：

\begin{equation}
\sum_{i=1}^n x_i \equiv x_1 + x_2 +\cdots + x_n
\end{equation}



\end{document}
